\section{Detection results per poison rate}
\label{app:results}

\Cref{tab:results-rate-1,tab:results-rate-5,tab:results-rate-10} give the full detection results at 1\%, 5\% and 10\% poisoning, 1 row per backdoored model, grouped by dataset. ASR is the attack success rate of the model, the share of triggered non-target images classified as the target, and CA is its clean accuracy on the test set. The 2 blocks of detection columns are PSBD-TM (token masking at the attention input) and PSBD-RD (dropout on the residual stream after both adds). AUROC is the one-sided area under the ROC curve, never flipped, so a value below 0.5 means the statistic ordered poisoned and clean inputs the wrong way round on that model. TPR@FPR is the share of triggered inputs flagged when the threshold is set on the 2000 clean validation images so that 10\% or 20\% of them are flagged.

A model whose attack success is below \AsrBar{} prints its ASR and CA but no detection numbers, since a detector cannot be credited or blamed for a backdoor that was never planted. Reading the 3 tables together shows what the evaluation can and cannot say. BadNets, Blend, BPP and LF implant at every rate on every dataset. TaCT implants everywhere except GTSRB at 1\%. WaNet implants only at 5\% and 10\% and not on CIFAR-100, SIG only on CIFAR-10 at 10\% and on SVHN and EuroSAT, and Adaptive-Blend never reaches the bar. \Cref{tab:results-benign} gives the clean accuracy of the benign reference model of each dataset, trained with the same recipe, which every clean-accuracy drop in the paper subtracts against.

% generated by scripts/paper/tab_results_by_rate.py from results/coverage/coverage.json, results/<folder>/psbd_metrics.json at 2026-09-23T18:28:30.800433+00:00, commit 5a1f1c5814dc5676c093f3021490a849712a13f4-dirty. Do not edit by hand.
\begin{table*}[htbp]
\centering
\small
\setlength{\tabcolsep}{4pt}
\caption{Detection results at 1\% poisoning on ViT-B/16.}
\label{tab:results-rate-1}
\begin{tabular}{lrrrrrrrr}
\toprule
 & & & \multicolumn{3}{c}{PSBD-TM} & \multicolumn{3}{c}{PSBD-RD} \\
\cmidrule(lr){4-6} \cmidrule(lr){7-9}
Attack & ASR & CA & AUROC & \multicolumn{2}{c}{TPR@FPR} & AUROC & \multicolumn{2}{c}{TPR@FPR} \\
\cmidrule(lr){5-6} \cmidrule(lr){8-9}
 & & & & 10\% & 20\% & & 10\% & 20\% \\
\midrule
\multicolumn{9}{l}{\textit{CIFAR-10}} \\
BadNets & 0.997 & 0.952 & \textbf{0.982} & 0.959 & 0.978 & 0.309 & 0.015 & 0.046 \\
Blend & 1.000 & 0.954 & 0.922 & 0.773 & 0.869 & \textbf{0.990} & 0.983 & 0.992 \\
LF & 0.982 & 0.953 & \textbf{0.979} & 0.954 & 0.975 & 0.941 & 0.753 & 0.898 \\
BPP & 0.982 & 0.952 & \textbf{0.986} & 0.956 & 0.976 & 0.970 & 0.897 & 0.952 \\
WaNet & 0.112 & 0.954 & -- & -- & -- & -- & -- & -- \\
TaCT & 0.985 & 0.938 & \textbf{0.979} & 0.005 & 0.009 & 0.464 & 0.013 & 0.038 \\
SIG & 0.340 & 0.951 & -- & -- & -- & -- & -- & -- \\
Adaptive-Blend & 0.622 & 0.953 & -- & -- & -- & -- & -- & -- \\
\midrule
\multicolumn{9}{l}{\textit{CIFAR-100}} \\
BadNets & 1.000 & 0.825 & \textbf{0.988} & 0.999 & 1.000 & 0.744 & 0.244 & 0.476 \\
Blend & 0.989 & 0.829 & \textbf{0.976} & 0.959 & 0.983 & 0.792 & 0.347 & 0.598 \\
LF & 0.942 & 0.830 & 0.956 & 0.942 & 0.947 & \textbf{0.956} & 0.903 & 0.940 \\
BPP & 0.914 & 0.824 & 0.921 & 0.849 & 0.919 & \textbf{0.950} & 0.914 & 0.929 \\
WaNet & 0.057 & 0.822 & -- & -- & -- & -- & -- & -- \\
TaCT & 0.989 & 0.820 & \textbf{0.909} & 0.506 & 0.690 & 0.493 & 0.023 & 0.103 \\
SIG & 0.250 & 0.817 & -- & -- & -- & -- & -- & -- \\
Adaptive-Blend & 0.536 & 0.823 & -- & -- & -- & -- & -- & -- \\
\midrule
\multicolumn{9}{l}{\textit{GTSRB}} \\
BadNets & 1.000 & 0.994 & \textbf{0.999} & 0.996 & 0.999 & 0.222 & 0.001 & 0.008 \\
Blend & 1.000 & 0.991 & \textbf{1.000} & 1.000 & 1.000 & 0.999 & 1.000 & 1.000 \\
LF & 0.979 & 0.990 & 0.976 & 0.926 & 0.973 & \textbf{0.987} & 0.977 & 0.985 \\
BPP & 0.868 & 0.986 & \textbf{0.898} & 0.717 & 0.809 & 0.755 & 0.128 & 0.507 \\
WaNet & 0.119 & 0.979 & -- & -- & -- & -- & -- & -- \\
TaCT & 0.000 & 0.121 & -- & -- & -- & -- & -- & -- \\
SIG & 0.356 & 0.981 & -- & -- & -- & -- & -- & -- \\
Adaptive-Blend & 0.560 & 0.983 & -- & -- & -- & -- & -- & -- \\
\midrule
\multicolumn{9}{l}{\textit{Tiny ImageNet}} \\
BadNets & 0.999 & 0.752 & \textbf{0.980} & 0.989 & 0.997 & 0.892 & 0.700 & 0.836 \\
Blend & 0.999 & 0.760 & 0.992 & 0.985 & 0.992 & \textbf{0.996} & 0.994 & 0.998 \\
LF & 0.922 & 0.762 & 0.941 & 0.929 & 0.935 & \textbf{0.952} & 0.918 & 0.934 \\
BPP & 0.971 & 0.748 & 0.980 & 0.972 & 0.975 & \textbf{0.981} & 0.972 & 0.974 \\
WaNet & 0.178 & 0.747 & -- & -- & -- & -- & -- & -- \\
TaCT & 0.976 & 0.750 & \textbf{0.575} & 0.024 & 0.119 & 0.487 & 0.024 & 0.071 \\
SIG & 0.076 & 0.762 & -- & -- & -- & -- & -- & -- \\
Adaptive-Blend & 0.554 & 0.747 & -- & -- & -- & -- & -- & -- \\
\midrule
\multicolumn{9}{l}{\textit{SVHN}} \\
SIG & 0.700 & 0.960 & -- & -- & -- & -- & -- & -- \\
\midrule
\multicolumn{9}{l}{\textit{EuroSAT}} \\
SIG & 0.813 & 0.969 & -- & -- & -- & -- & -- & -- \\
\bottomrule
\end{tabular}
\end{table*}

% generated by scripts/paper/tab_results_by_rate.py from results/coverage/coverage.json, results/<folder>/psbd_metrics.json at 2026-09-23T16:16:21.848434+00:00, commit afe5fd94d50b77f30f654fe075b768040987adef-dirty. Do not edit by hand.
\begin{table*}[htbp]
\centering
\small
\setlength{\tabcolsep}{4pt}
\caption{Detection results at 5\% poisoning on ViT-B/16.}
\label{tab:results-rate-5}
\begin{tabular}{lrrrrrrrr}
\toprule
 & & & \multicolumn{3}{c}{PSBD-TM} & \multicolumn{3}{c}{PSBD-RD} \\
\cmidrule(lr){4-6} \cmidrule(lr){7-9}
Attack & ASR & CA & AUROC & \multicolumn{2}{c}{TPR@FPR} & AUROC & \multicolumn{2}{c}{TPR@FPR} \\
\cmidrule(lr){5-6} \cmidrule(lr){8-9}
 & & & & 10\% & 20\% & & 10\% & 20\% \\
\midrule
\multicolumn{9}{l}{\textit{CIFAR-10}} \\
BadNets & 1.000 & 0.945 & \textbf{0.992} & 0.965 & 0.995 & 0.630 & 0.031 & 0.168 \\
Blend & 1.000 & 0.957 & 0.971 & 0.908 & 0.989 & \textbf{0.984} & 0.973 & 0.995 \\
LF & 0.997 & 0.944 & \textbf{0.991} & 0.983 & 0.994 & 0.959 & 0.837 & 0.947 \\
BPP & 0.987 & 0.950 & 0.801 & 0.587 & 0.655 & \textbf{0.880} & 0.673 & 0.788 \\
WaNet & 0.961 & 0.913 & 0.937 & 0.838 & 0.913 & \textbf{0.956} & 0.908 & 0.965 \\
TaCT & 0.996 & 0.946 & \textbf{0.966} & 0.000 & 0.031 & 0.834 & 0.006 & 0.020 \\
SIG & 0.599 & 0.953 & -- & -- & -- & -- & -- & -- \\
Adaptive-Blend & 0.594 & 0.946 & -- & -- & -- & -- & -- & -- \\
\midrule
\multicolumn{9}{l}{\textit{CIFAR-100}} \\
BadNets & 1.000 & 0.825 & \textbf{0.994} & 0.999 & 1.000 & 0.826 & 0.335 & 0.672 \\
Blend & 1.000 & 0.830 & 0.979 & 0.941 & 0.976 & \textbf{0.999} & 0.999 & 1.000 \\
LF & 0.953 & 0.829 & \textbf{0.967} & 0.954 & 0.959 & 0.966 & 0.934 & 0.955 \\
BPP & 0.988 & 0.833 & 0.985 & 0.961 & 0.977 & \textbf{0.988} & 0.973 & 0.984 \\
WaNet & 0.643 & 0.809 & -- & -- & -- & -- & -- & -- \\
TaCT & 1.000 & 0.825 & \textbf{0.933} & 0.540 & 0.862 & 0.554 & 0.218 & 0.552 \\
SIG & 0.249 & 0.812 & -- & -- & -- & -- & -- & -- \\
Adaptive-Blend & 0.579 & 0.812 & -- & -- & -- & -- & -- & -- \\
\midrule
\multicolumn{9}{l}{\textit{GTSRB}} \\
BadNets & 0.573 & 0.086 & -- & -- & -- & -- & -- & -- \\
Blend & 0.951 & 0.080 & -- & -- & -- & -- & -- & -- \\
LF & 0.999 & 0.974 & \textbf{0.998} & 0.999 & 0.999 & 0.997 & 0.997 & 0.999 \\
BPP & 0.991 & 0.978 & 0.992 & 0.978 & 0.987 & \textbf{0.994} & 0.988 & 0.992 \\
WaNet & 0.830 & 0.971 & -- & -- & -- & -- & -- & -- \\
TaCT & 1.000 & 0.986 & \textbf{0.942} & 0.502 & 0.818 & 0.392 & 0.010 & 0.041 \\
SIG & 0.673 & 0.990 & -- & -- & -- & -- & -- & -- \\
Adaptive-Blend & 0.777 & 0.981 & -- & -- & -- & -- & -- & -- \\
\midrule
\multicolumn{9}{l}{\textit{Tiny ImageNet}} \\
BadNets & 1.000 & 0.752 & \textbf{0.984} & 0.999 & 1.000 & 0.837 & 0.566 & 0.739 \\
Blend & 0.999 & 0.754 & 0.997 & 0.997 & 0.999 & \textbf{0.999} & 0.999 & 0.999 \\
LF & 0.987 & 0.753 & 0.987 & 0.987 & 0.988 & \textbf{0.991} & 0.987 & 0.989 \\
BPP & 0.985 & 0.746 & 0.987 & 0.984 & 0.986 & \textbf{0.987} & 0.977 & 0.987 \\
WaNet & 0.922 & 0.748 & 0.930 & 0.924 & 0.928 & \textbf{0.953} & 0.922 & 0.938 \\
TaCT & 0.952 & 0.743 & \textbf{0.663} & 0.381 & 0.500 & 0.488 & 0.024 & 0.119 \\
SIG & 0.432 & 0.748 & -- & -- & -- & -- & -- & -- \\
Adaptive-Blend & 0.783 & 0.745 & -- & -- & -- & -- & -- & -- \\
\midrule
\multicolumn{9}{l}{\textit{SVHN}} \\
SIG & 0.880 & 0.958 & \textbf{0.797} & 0.120 & 0.482 & 0.646 & 0.209 & 0.339 \\
\midrule
\multicolumn{9}{l}{\textit{EuroSAT}} \\
SIG & 0.873 & 0.979 & pending & pending & pending & 0.593 & 0.143 & 0.276 \\
\bottomrule
\end{tabular}
\end{table*}

% generated by scripts/paper/tab_results_by_rate.py from results/coverage/coverage.json, results/<folder>/psbd_metrics.json at 2026-09-23T14:43:32.345383+00:00, commit 039fc6fea9212f8f6d49218ed14b336aa8f6a0d3-dirty. Do not edit by hand.
\begin{table*}[htbp]
\centering
\small
\setlength{\tabcolsep}{4pt}
\caption{Detection results at 10\% poisoning on ViT-B/16.}
\label{tab:results-rate-10}
\begin{tabular}{lrrrrrrrr}
\toprule
 & & & \multicolumn{3}{c}{PSBD-TM} & \multicolumn{3}{c}{PSBD-RD} \\
\cmidrule(lr){4-6} \cmidrule(lr){7-9}
Attack & ASR & CA & AUROC & \multicolumn{2}{c}{TPR@FPR} & AUROC & \multicolumn{2}{c}{TPR@FPR} \\
\cmidrule(lr){5-6} \cmidrule(lr){8-9}
 & & & & 10\% & 20\% & & 10\% & 20\% \\
\midrule
\multicolumn{9}{l}{\textit{CIFAR-10}} \\
BadNets & 1.000 & 0.942 & \textbf{0.991} & 0.935 & 1.000 & 0.720 & 0.113 & 0.301 \\
Blend & 1.000 & 0.951 & 0.906 & 0.754 & 0.849 & \textbf{0.991} & 0.980 & 0.992 \\
LF & 0.999 & 0.951 & \textbf{0.993} & 0.995 & 0.999 & 0.977 & 0.908 & 0.972 \\
BPP & 0.994 & 0.954 & 0.871 & 0.736 & 0.830 & \textbf{0.875} & 0.720 & 0.799 \\
WaNet & 0.890 & 0.946 & 0.459 & 0.087 & 0.268 & \textbf{0.947} & 0.846 & 0.914 \\
TaCT & 1.000 & 0.858 & \textbf{0.963} & 0.030 & 0.206 & 0.261 & 0.000 & 0.000 \\
SIG & 0.901 & 0.846 & 0.418 & 0.071 & 0.097 & \textbf{0.919} & 0.810 & 0.878 \\
Adaptive-Blend & 0.622 & 0.951 & -- & -- & -- & -- & -- & -- \\
\midrule
\multicolumn{9}{l}{\textit{CIFAR-100}} \\
BadNets & 1.000 & 0.821 & \textbf{0.998} & 1.000 & 1.000 & 0.497 & 0.114 & 0.234 \\
Blend & 1.000 & 0.816 & 0.998 & 1.000 & 1.000 & \textbf{1.000} & 1.000 & 1.000 \\
LF & 0.995 & 0.823 & \textbf{0.995} & 0.995 & 0.995 & 0.982 & 0.959 & 0.983 \\
BPP & 0.991 & 0.817 & 0.985 & 0.963 & 0.975 & \textbf{0.985} & 0.960 & 0.979 \\
WaNet & 0.793 & 0.818 & -- & -- & -- & -- & -- & -- \\
TaCT & 0.989 & 0.821 & \textbf{0.891} & 0.161 & 0.586 & 0.488 & 0.000 & 0.103 \\
SIG & 0.178 & 0.820 & -- & -- & -- & -- & -- & -- \\
Adaptive-Blend & 0.604 & 0.808 & -- & -- & -- & -- & -- & -- \\
\midrule
\multicolumn{9}{l}{\textit{GTSRB}} \\
BadNets & 1.000 & 0.972 & \textbf{1.000} & 1.000 & 1.000 & 0.999 & 0.999 & 1.000 \\
Blend & 1.000 & 0.988 & \textbf{1.000} & 1.000 & 1.000 & 0.998 & 0.996 & 0.998 \\
LF & 0.999 & 0.988 & 0.998 & 0.998 & 0.999 & \textbf{0.999} & 0.999 & 0.999 \\
BPP & 0.995 & 0.986 & \textbf{0.989} & 0.971 & 0.986 & 0.977 & 0.971 & 0.992 \\
WaNet & 0.947 & 0.958 & 0.948 & 0.949 & 0.952 & \textbf{0.960} & 0.950 & 0.960 \\
TaCT & 1.000 & 0.931 & \textbf{0.748} & 0.213 & 0.843 & 0.295 & 0.000 & 0.000 \\
Adaptive-Blend & 0.838 & 0.976 & -- & -- & -- & -- & -- & -- \\
\midrule
\multicolumn{9}{l}{\textit{Tiny ImageNet}} \\
BadNets & 1.000 & 0.752 & \textbf{0.993} & 1.000 & 1.000 & 0.973 & 0.951 & 0.985 \\
Blend & 1.000 & 0.749 & 0.998 & 0.999 & 0.999 & \textbf{0.999} & 0.999 & 1.000 \\
LF & 0.973 & 0.740 & 0.973 & 0.973 & 0.977 & \textbf{0.982} & 0.970 & 0.976 \\
BPP & 0.996 & 0.757 & 0.984 & 0.981 & 0.991 & \textbf{0.997} & 0.995 & 0.996 \\
WaNet & 0.968 & 0.747 & 0.950 & 0.932 & 0.970 & \textbf{0.966} & 0.924 & 0.963 \\
TaCT & 0.929 & 0.756 & \textbf{0.810} & 0.381 & 0.619 & 0.417 & 0.048 & 0.071 \\
SIG & 0.122 & 0.747 & -- & -- & -- & -- & -- & -- \\
Adaptive-Blend & 0.505 & 0.753 & -- & -- & -- & -- & -- & -- \\
\midrule
\multicolumn{9}{l}{\textit{SVHN}} \\
Blend & 1.000 & 0.951 & \textbf{0.941} & 0.731 & 0.890 & 0.623 & 0.188 & 0.358 \\
SIG & 0.971 & 0.958 & pending & pending & pending & 0.689 & 0.203 & 0.367 \\
\midrule
\multicolumn{9}{l}{\textit{EuroSAT}} \\
Blend & 1.000 & 0.980 & 0.846 & 0.014 & 0.251 & \textbf{0.894} & 0.501 & 0.739 \\
SIG & 0.922 & 0.972 & \textbf{0.617} & 0.018 & 0.135 & 0.559 & 0.094 & 0.206 \\
\bottomrule
\end{tabular}
\end{table*}

% generated by scripts/paper/tab_results_by_rate.py from results/coverage/coverage.json at 2026-09-23T14:43:39.078303+00:00, commit 039fc6fea9212f8f6d49218ed14b336aa8f6a0d3-dirty. Do not edit by hand.
\begin{table}[htbp]
\centering
\small
\caption{Benign reference models.}
\label{tab:results-benign}
\begin{adjustbox}{max width=\linewidth}
\begin{tabular}{lrr}
\toprule
Dataset & Clean Accuracy & Attacks Above ASR Bar \\
\midrule
CIFAR-10 & 0.953 & 18 \\
CIFAR-100 & 0.811 & 15 \\
GTSRB & 0.991 & 15 \\
Tiny ImageNet & 0.755 & 17 \\
SVHN & 0.964 & 3 \\
EuroSAT & 0.977 & 3 \\
\bottomrule
\end{tabular}
\end{adjustbox}
\end{table}

\FloatBarrier

\input{sections/attacks}
\FloatBarrier

\section{The placement basis}
\label{app:basis}

\Cref{tab:basis-declaration} lists the placements every model was swept over. A placement is a site (where the perturbation is injected, \cref{fig:vit-block}) paired with a perturbation (what the hook does), optionally restricted to a band of blocks. The family column groups sites by where they sit relative to the residual stream. Input-side sites perturb a sublayer's input before its LayerNorm. Residual-adjacent sites perturb a branch output just before its residual add or the stream just after it. The structured sites perturb attention heads and MLP hidden units, and the ported sites come from other detectors, scaling input pixels as SCALE-UP does and scaling a LayerNorm's gain as IBD-PSC does. The rates column is the ladder each placement was swept over. Masks and dropout use a share of the tensor, Gaussian noise a multiple of the tensor's standard deviation. The last entries were added after the first had been read, to fill \cref{tab:staircase-residual,tab:staircase-dropout-sites,tab:staircase-bands}. All \BasisSize{} declared placements now reach the adaptive shift target somewhere on the panel. Coverage is still uneven across them, so every claim about the basis is restricted to the placements common to every compared cell.

% generated by scripts/paper/app_basis.py from configs/psbd_basis.json at 2026-09-10T23:32:03.398466+00:00, commit 831df7ef5e8bc9285264b261f848e1f2cad5ef61-dirty. Do not edit by hand.
\begin{table}[t]
\centering
\small
\caption{The placement basis as configs/psbd\_basis.json declares it: every panel cell carries every 1 of these placements, each swept over the rate ladder in the last column.}
\label{tab:basis-declaration}
\begin{tabular}{llllll}
\toprule
placement id & position & operator & blocks & family & rates \\
\midrule
before\_attention\_norm\_token\_mask & before\_attention\_norm & token\_mask & all & input\_side & 10 (0.05 to 0.9) \\
before\_attention\_norm\_channel\_mask & before\_attention\_norm & channel\_mask & all & input\_side & 10 (0.05 to 0.9) \\
before\_attention\_norm\_gaussian & before\_attention\_norm & gaussian & all & input\_side & 12 (0.05 to 8) \\
before\_mlp\_norm\_token\_mask & before\_mlp\_norm & token\_mask & all & input\_side & 10 (0.05 to 0.9) \\
both\_sublayer\_inputs\_token\_mask & both\_sublayer\_inputs & token\_mask & all & input\_side & 9 (0.1 to 0.9) \\
input\_pixels\_scale\_up & input\_pixels & scale\_up & all & input\_side & 11 (0.5 to 25) \\
before\_mlp\_gaussian & before\_mlp & gaussian & all & input\_side & 12 (0.05 to 8) \\
before\_attention\_residual\_token\_mask & before\_attention\_residual & token\_mask & all & residual\_adjacent & 10 (0.05 to 0.9) \\
after\_attention\_residual\_token\_mask & after\_attention\_residual & token\_mask & all & residual\_adjacent & 10 (0.05 to 0.9) \\
pre\_residual & pre\_residual & dropout & all & residual\_adjacent & 9 (0.1 to 0.9) \\
post\_residual & post\_residual & dropout & all & residual\_adjacent & 16 (0.005 to 0.9) \\
mlp\_neurons\_channel\_mask & mlp\_neurons & channel\_mask & all & structured & 10 (0.05 to 0.9) \\
mlp\_norm\_out\_gain\_scale & mlp\_norm\_out & gain\_scale & all & ported & 11 (0.25 to 15) \\
before\_attention\_norm\_blocks\_5\_8\_token\_mask & before\_attention\_norm & token\_mask & 5-8 & depth\_band & 12 (0.05 to 0.99) \\
before\_attention\_norm\_blocks\_9\_12\_token\_mask & before\_attention\_norm & token\_mask & 9-12 & depth\_band & 12 (0.05 to 0.99) \\
pre\_residual\_blocks\_1\_4 & pre\_residual & dropout & 1-4 & depth\_band & 11 (0.1 to 0.99) \\
pre\_residual\_blocks\_5\_8 & pre\_residual & dropout & 5-8 & depth\_band & 11 (0.1 to 0.99) \\
pre\_residual\_blocks\_9\_12 & pre\_residual & dropout & 9-12 & depth\_band & 11 (0.1 to 0.99) \\
\bottomrule
\end{tabular}
\end{table}

\FloatBarrier

\Cref{tab:basis-ranking} ranks the placements with a complete sweep by mean AUROC, and gives each placement's lowest single-model AUROC. The top 3 placements all mask whole tokens on an input to attention. The gain-scaling and noise placements sit in the bottom half, and scaling input pixels is last. PSBD-TM and token masking on the attention branch before its residual add share the top of the ranking at the same mean AUROC to 3 decimal places. \Cref{app:protocol} says why we recommend the attention input. \Cref{fig:basis} plots this table.

% generated by scripts/paper/app_basis.py from results/coverage/coverage.json, configs/psbd_basis.json, results/<folder>/psbd_metrics.json (71 cells) at 2026-09-23T16:00:54.049534+00:00, commit afe5fd94d50b77f30f654fe075b768040987adef-dirty. Do not edit by hand.
\begin{table}[htbp]
\centering
\small
\caption{Every placement of the basis ranked by mean AUROC on ViT-B/16. n is the number of models whose rate ladder reaches the shift ratio target, and the last column is the lowest single-model AUROC.}
\label{tab:basis-ranking}
\begin{adjustbox}{max width=\linewidth}
\begin{tabular}{rllrrr}
\toprule
Rank & Placement & Family & n & AUROC & Lowest AUROC \\
\midrule
1 & token mask, attention input & input side & 69 & 0.927 & 0.418 \\
2 & token mask, attention output before the add & residual adjacent & 68 & 0.927 & 0.537 \\
3 & token mask, attention input, blocks 1 to 4 & depth band & 12 & 0.924 & 0.653 \\
4 & token mask, both sublayer inputs & input side & 71 & 0.898 & 0.466 \\
5 & gain scale, MLP norm output & ported & 70 & 0.896 & 0.266 \\
6 & token mask, attention input, blocks 9 to 12 & depth band & 31 & 0.893 & 0.514 \\
7 & token mask, attention input, blocks 5 to 8 & depth band & 44 & 0.883 & 0.244 \\
8 & channel mask, attention input & input side & 71 & 0.883 & 0.091 \\
9 & dropout, before both residual adds, blocks 9 to 12 & depth band & 46 & 0.880 & 0.372 \\
10 & dropout, before both residual adds, blocks 5 to 8 & depth band & 71 & 0.873 & 0.365 \\
11 & noise, MLP input after norm & input side & 71 & 0.871 & 0.226 \\
12 & dropout, attention input & input side & 65 & 0.859 & 0.074 \\
13 & dropout, before both residual adds & residual adjacent & 71 & 0.838 & 0.093 \\
14 & token mask, MLP input & input side & 71 & 0.826 & 0.048 \\
15 & dropout, after both residual adds & residual adjacent & 71 & 0.818 & 0.222 \\
16 & dropout, attention input after norm & input side & 65 & 0.811 & 0.165 \\
17 & dropout, embedding output & input side & 65 & 0.810 & 0.172 \\
18 & channel mask, MLP neurons & structured & 71 & 0.808 & 0.250 \\
19 & noise, attention output before the add & residual adjacent & 65 & 0.806 & 0.055 \\
20 & dropout, before both residual adds, blocks 1 to 4 & depth band & 71 & 0.805 & 0.080 \\
21 & dropout, stream after the attention add & residual adjacent & 65 & 0.804 & 0.190 \\
22 & noise, attention input & input side & 71 & 0.804 & 0.031 \\
23 & channel mask, attention output before the add & residual adjacent & 65 & 0.789 & 0.029 \\
24 & token mask, stream after the attention add & residual adjacent & 71 & 0.761 & 0.050 \\
25 & dropout, MLP input & input side & 65 & 0.759 & 0.070 \\
26 & scale up, input pixels & input side & 71 & 0.756 & 0.225 \\
27 & dropout, attention output before the add & residual adjacent & 0 & -- & -- \\
\bottomrule
\end{tabular}
\end{adjustbox}
\end{table}

\FloatBarrier

\section{Depth bands at matched disturbance}
\label{app:bands}

\Cref{tab:depth-bands} reads every depth band at the matched rule, a clean shift ratio of \PlacementMatchTarget{}, where more of the bands are read on the same models than at the adaptive target. The adaptive-rule readings in \cref{tab:staircase-bands} put the early input-side band near the top on the \StaircaseBandsBeforeAttentionNormBlocksOneFourTokenMaskN{} models that reach the target there, and at matched disturbance on \BandOneFourMinusAllInputSideN{} models it is the worst band of the 3. A band's n is worth reading beside its mean.

% generated by scripts/paper/tab_depth_bands.py from results/coverage/coverage.json, results/<folder>/psbd_metrics.json (65 cells) at 2026-09-11T04:12:43.663538+00:00, commit 4a04e5dd204b14d8368fb7eeef7092bd331e00d9-dirty. Do not edit by hand.
\begin{table}[htbp]
\centering
\small
\caption{Depth-band placements at the matched 0.6 rate, mean AUROC at q0.25, the paired delta against the family's all-blocks placement with a 5000-resample bootstrap 95\% interval, and the mean achieved clean-validation shift ratio at the chosen rate. Token masking at the attention input restricted to blocks 1 to 4 has not been swept yet on any model in the panel of 65 backdoored models and prints pending throughout.}
\label{tab:depth-bands}
\begin{adjustbox}{max width=\linewidth}
\begin{tabular}{lrrrlr}
\toprule
placement & n & mean AUROC matched06 & delta vs all blocks & 95\% CI & mean achieved shift \\
\midrule
input-side, all blocks & 65 & 0.926 & -- & -- & 0.595 \\
input-side, blocks 5-8 & 65 & 0.863 & -0.063 & [-0.100, -0.031] & 0.603 \\
input-side, blocks 9-12 & 65 & 0.864 & -0.062 & [-0.108, -0.021] & 0.489 \\
input-side, blocks 1-4 & \pending & \pending & \pending & \pending & \pending \\
residual-adjacent, all blocks & 65 & 0.841 & -- & -- & 0.658 \\
residual-adjacent, blocks 1-4 & 65 & 0.779 & -0.062 & [-0.081, -0.042] & 0.582 \\
residual-adjacent, blocks 5-8 & 65 & 0.893 & +0.052 & [+0.027, +0.079] & 0.565 \\
residual-adjacent, blocks 9-12 & 65 & 0.874 & +0.033 & [-0.012, +0.076] & 0.598 \\
\bottomrule
\end{tabular}
\end{adjustbox}
\end{table}

\FloatBarrier

\section{Every attention-side probe}
\label{app:attention}

Attention is the only operation in a ViT block that moves information between tokens, so the mechanism of \cref{sec:mechanism} makes it the interesting sublayer to perturb. We measured \AttentionProbesTotal{} probes on it, \AttentionProbesInBasis{} of them declared by the basis and \AttentionProbesOutsideBasis{} not.

\Cref{tab:attention-probes} gives them all. The probes the basis does not declare cannot join the ranking of \cref{tab:basis-ranking}, because a placement added to a ranking after its score is read is the oracle selection the declaration exists to prevent, so they are reported separately with their own coverage. The point of reporting them is that \AttentionProbesBeatingRecommended{} of them beats the placement we recommend. The best is \AttentionProbesBestOutsideName{} at \AttentionProbesBestOutside{} on \AttentionProbesBestOutsideN{} models, below the recommended placement and on far fewer models. Masking whole attention heads, the most attention-specific probe available, reads \AttentionHeadMaskAuroc{} on \AttentionHeadMaskN{} models with \AttentionHeadMaskTprTen{} TPR at a 10\% budget, which places it among the stronger probes without displacing token masking at the attention input.

% generated by scripts/paper/tab_attention.py from results/coverage/coverage.json (71 clearing cells), results/*/psbd_metrics.json at 2026-09-23T16:56:31.219165+00:00, commit d4a8014354b8fcc8d014beec50fef60ed46fa47d-dirty. Do not edit by hand.
\begin{table}[htbp]
\centering
\small
\caption{Every attention-side probe with at least 5 backdoored ViT-B/16 models, at the adaptive rule and the headline quantile. The second column says whether the pre-registered basis declares the placement. The probes it does not declare are not part of the ranking in \cref{tab:basis-ranking} and were not eligible to be, since adding a placement to a ranking after reading its score is what the declaration exists to prevent. Coverage is uneven, so rows are not paired.}
\label{tab:attention-probes}
\begin{adjustbox}{max width=\linewidth}
\begin{tabular}{llrrrr}
\toprule
probe & in basis & n & AUROC & TPR@10 & TPR@20 \\
\midrule
token mask, attention input & yes & 69 & 0.927 & 0.77 & 0.83 \\
token mask, attention output before the add & yes & 68 & 0.927 & 0.79 & 0.88 \\
noise, attention input after norm & no & 19 & 0.913 & 0.79 & 0.87 \\
head mask, attention heads & no & 37 & 0.896 & 0.74 & 0.84 \\
channel mask, attention input & yes & 71 & 0.883 & 0.68 & 0.77 \\
token mask, attention input after norm & no & 37 & 0.876 & 0.75 & 0.81 \\
drop path, attention output before the add & no & 37 & 0.872 & 0.66 & 0.79 \\
dropout, attention input & yes & 65 & 0.859 & 0.63 & 0.73 \\
channel mask, attention input after norm & no & 37 & 0.847 & 0.66 & 0.77 \\
channel mask, stream after the attention add & no & 37 & 0.835 & 0.65 & 0.74 \\
noise, stream after the attention add & no & 37 & 0.820 & 0.59 & 0.70 \\
dropout, attention output before the add & yes & 65 & 0.813 & 0.59 & 0.69 \\
dropout, attention input after norm & yes & 65 & 0.811 & 0.58 & 0.70 \\
noise, attention output before the add & yes & 65 & 0.806 & 0.60 & 0.68 \\
dropout, stream after the attention add & yes & 65 & 0.804 & 0.60 & 0.69 \\
noise, attention input & yes & 71 & 0.804 & 0.54 & 0.66 \\
channel mask, attention output before the add & yes & 65 & 0.789 & 0.56 & 0.66 \\
token mask, stream after the attention add & yes & 71 & 0.761 & 0.43 & 0.56 \\
gain scale, attention norm output & no & 37 & 0.639 & 0.39 & 0.47 \\
\bottomrule
\end{tabular}
\end{adjustbox}
\end{table}

\FloatBarrier

\section{Whether placements agree on which inputs are fragile}
\label{app:agreement}

If every placement read the same per-sample quantity, a union of them could add nothing and the family grouping would sort placements that behave alike. \Cref{tab:mech-operator-agreement,tab:mech-operator-agreement-family} test both. Each placement is read at its own matched rate, so the comparison is at equal disturbance, and the correlation is Spearman's over the clean split of every model carrying both placements.

% generated by scripts/paper/mech_operator_agreement.py from results/<folder>/psbd/<placement>/rate_*_{clean,backdoor}.pt at 2026-09-23T19:56:39.614936+00:00, commit 4087ff7b8dddaf7e969f38f703166a0e490aaacf-dirty. Do not edit by hand.
\begin{table}[htbp]
\centering
\small
\caption{Whether placements agree on which inputs are fragile. Mean Spearman correlation, over the models whose attack succeeded, between each basis placement's per-sample fractional PSU and that of token masking at the attention input, each at its own matched 0.6 rate. n is the number of models carrying both placements at a usable rate.}
\label{tab:mech-operator-agreement}
\begin{adjustbox}{max width=\linewidth}
\begin{tabular}{lrrrrr}
\toprule
placement & family & rho vs token mask, attention input (clean) & n & rho vs token mask, attention input (backdoor) & n \\
\midrule
channel mask, attention input & input side & 0.642 & 71 & 0.311 & 71 \\
noise, attention input & input side & 0.396 & 71 & 0.273 & 71 \\
token mask, MLP input & input side & 0.581 & 71 & 0.470 & 71 \\
token mask, both sublayer inputs & input side & 0.806 & 71 & 0.607 & 71 \\
scale up, input pixels & input side & 0.203 & 71 & 0.125 & 71 \\
noise, MLP input after norm & input side & 0.553 & 71 & 0.303 & 71 \\
token mask, attention output before the add & residual adjacent & 0.538 & 71 & 0.457 & 71 \\
token mask, stream after the attention add & residual adjacent & 0.573 & 71 & 0.477 & 71 \\
dropout, stream after the attention add & residual adjacent & 0.531 & 65 & 0.304 & 65 \\
dropout, before both residual adds & residual adjacent & 0.606 & 71 & 0.326 & 71 \\
dropout, after both residual adds & residual adjacent & 0.611 & 71 & 0.358 & 71 \\
channel mask, MLP neurons & structured & 0.508 & 71 & 0.225 & 71 \\
gain scale, MLP norm output & ported & 0.384 & 71 & 0.322 & 71 \\
token mask, attention input, blocks 5 to 8 & depth band & 0.565 & 71 & 0.451 & 71 \\
token mask, attention input, blocks 9 to 12 & depth band & 0.684 & 71 & 0.192 & 71 \\
dropout, before both residual adds, blocks 1 to 4 & depth band & 0.564 & 71 & 0.286 & 71 \\
dropout, before both residual adds, blocks 5 to 8 & depth band & 0.654 & 71 & 0.319 & 71 \\
dropout, before both residual adds, blocks 9 to 12 & depth band & 0.438 & 71 & 0.247 & 71 \\
dropout, attention input & input side & 0.658 & 65 & 0.308 & 65 \\
dropout, attention input after norm & input side & 0.403 & 65 & 0.180 & 65 \\
dropout, MLP input & input side & 0.482 & 65 & 0.217 & 65 \\
dropout, embedding output & input side & 0.390 & 65 & 0.202 & 65 \\
token mask, attention input, blocks 1 to 4 & depth band & 0.514 & 65 & 0.343 & 65 \\
dropout, attention output before the add & residual adjacent & 0.540 & 65 & 0.213 & 65 \\
channel mask, attention output before the add & residual adjacent & 0.517 & 65 & 0.218 & 65 \\
noise, attention output before the add & residual adjacent & 0.572 & 65 & 0.307 & 65 \\
\bottomrule
\end{tabular}
\end{adjustbox}
\end{table}

% generated by scripts/paper/mech_operator_agreement.py from results/<folder>/psbd/<placement>/rate_*_{clean,backdoor}.pt at 2026-09-23T18:19:00.094983+00:00, commit 5a1f1c5814dc5676c093f3021490a849712a13f4-dirty. Do not edit by hand.
\begin{table}[htbp]
\centering
\small
\caption{B1: mean per-sample PSU agreement (Spearman rho, averaged over placement pairs and over models) within and across the input-side and residual-adjacent families, the axis H20 splits on.}
\label{tab:mech-operator-agreement-family}
\begin{adjustbox}{max width=\linewidth}
\begin{tabular}{lrrrr}
\toprule
comparison & clean rho & n pairs & backdoor rho & n pairs \\
\midrule
within input\_side & 0.472 & 55 & 0.264 & 55 \\
within residual\_adjacent & 0.582 & 28 & 0.349 & 28 \\
across families (input\_side vs residual\_adjacent) & 0.512 & 88 & 0.302 & 88 \\
\bottomrule
\end{tabular}
\end{adjustbox}
\end{table}

\FloatBarrier

\section{Selection protocol}
\label{app:protocol}

Choosing the best placement on the models it is later reported on would give an optimistic reading, so the basis declares a split before any sweep. Placements are selected on CIFAR-10 and GTSRB and reported on CIFAR-100 and Tiny ImageNet.

Run on ViT alone, that protocol would not have picked the placement we recommend. On the \GainsSelectionHalfN{} selection-half models its twin, token masking on the attention branch output, reads \GainsSelectionHalfTwin{} against \GainsSelectionHalfOurs{}, a gap of \GainsSelectionHalfTwinGap{}. On the \GainsReportHalfN{} report-half models the order reverses, \GainsReportHalfOurs{} against \GainsReportHalfTwin{}. The 2 are tied over the whole panel, so the ViT protocol cannot separate them and the data it selects on leans to the twin. The recommendation rests on 3 things the ViT selection half does not supply: on Swin-S the same protocol picks the attention input by \SwinSelectionHalfTwinGap{} (\cref{sec:swin}), the attention input is the more stable of the 2 across training seeds, and it is the site whose mechanism \cref{sec:mechanism} can state. The numbers in the main text pool both halves.

\section{The shift ratio target}
\label{app:rate}

The original rule reads each placement at the rate where 80\% of clean validation predictions change. \Cref{fig:shift-ladder} checks that target on our models by interpolating every model's rate ladder onto a shared grid of achieved shift ratios from 0.2 to 0.98 and reading AUROC and TPR along it, for all models and separately for the models at each poison rate.

\begin{figure}[htbp]
\centering
\includegraphics[width=\textwidth]{fig_shift_ladder.pdf}
\caption{Mean AUROC and TPR at the 10\% and 20\% false-positive budgets of PSBD-TM against the achieved clean shift ratio. The dashed line marks the adaptive target of \AdaptiveShiftTarget{}. Above it fewer models reach the target within the swept rate ladder and the mean is taken over those that do.}
\label{fig:shift-ladder}
\end{figure}

AUROC barely moves with the target, at \LadderAurocAtZeroSix{} at a shift ratio of \PlacementMatchTarget{} and \LadderAurocAtZeroEight{} at \AdaptiveShiftTarget{}, while the true-positive rate at a fixed false-positive budget keeps rising with it. At a 10\% budget it goes from \LadderTprOneZeroAtZeroSix{} to \LadderTprOneZeroAtZeroEight{} and \LadderTprOneZeroAtZeroNine{} at a shift ratio of 0.9. A higher target therefore buys operating-point performance without moving the ranking statistic. The caveat is that only \LadderAurocAtZeroNineN{} of the \LadderAurocAtZeroEightN{} models reach 0.9 within the swept rate ladder, so the means past \AdaptiveShiftTarget{} are over a shrinking subset. We keep the original target, which sits at the knee of every curve. These readings interpolate each model's ladder to the exact target, which is why they differ slightly from the adaptive rule's, which takes the smallest swept rate past it.

\section{Forward passes}
\label{app:passes}

\Cref{fig:forward-passes} plots the effect of the number of passes $k$ behind \cref{sec:robustness:passes}. Over all models the first 3 passes are available from the sweeps, and on 6 models at 1\% poisoning (BadNets, Blend and LF on CIFAR-100 and Tiny ImageNet) we swept $k$ up to 20. The sweeps at $k$ in 1, 5, 10 and 20 over the BPP, WaNet and TaCT models at 1\% and 5\% poisoning on CIFAR-100 and Tiny ImageNet are running and will replace the 6-model curve.

\begin{figure}[htbp]
\centering
\includegraphics[width=\textwidth]{fig_forward_passes.pdf}
\caption{Mean AUROC, TPR at 10\% FPR and TPR at 20\% FPR of PSBD-TM against the number of forward passes $k$. The full curve is on 6 models at 1\% poisoning with a $k = 20$ sweep. The short curve is on all models, whose sweeps hold $k \le 3$.}
\label{fig:forward-passes}
\end{figure}

\section{Per-layer profiles of the backdoor direction}
\label{app:layers}

\Cref{tab:tac-layers} gives, for each of the \TacLayersCheckpoints{} models in \cref{fig:tac-layers}, read on \TacLayersSamples{} paired images per model, the block at which the backdoor direction's norm relative to the clean activation peaks, that peak value, the first block at which the relative norm reaches half its final value (the onset), the final-layer linear centered kernel alignment (CKA) \cite{kornblith2019cka} between clean and triggered activations (a CKA of 1 means the trigger leaves the representation unchanged), and the mean absolute trigger-activated change (TAC) per dimension at the final layer. The benign rows are the control. Their peak relative norm never exceeds \TacLayersBenignPeakMax{} and their CKA is at least \TacLayersBenignCkaMin{}, so the profiles of the backdoored rows are learned structure and not an artifact of the trigger's pixels. TaCT stands out with a peak of \TacLayersTactPeakMin{} to \TacLayersTactPeakMax{} and a CKA of \TacLayersTactCkaMin{} to \TacLayersTactCkaMax{}. A source-specific attack changes the representation of the few source classes only, so a direction averaged over all classes is small even though the attack succeeds.

% generated by scripts/paper/mech_tac_layers.py from results/_experiments/tac_layers/tac_layers.json at 2026-09-23T18:21:02.095010+00:00, commit 5a1f1c5814dc5676c093f3021490a849712a13f4-dirty. Do not edit by hand.
\begin{table}[htbp]
\centering
\small
\caption{Per-layer profile of the backdoor direction on 18 ViT-B/16 models, 500 paired images each.}
\label{tab:tac-layers}
\begin{adjustbox}{max width=\linewidth}
\begin{tabular}{llrrrrr}
\toprule
Dataset & Model & Peak layer & Peak rel. norm & Onset layer & Final CKA & Final TAC \\
\midrule
GTSRB & BadNets 10\% & 12 & 1.64 & 7 & 0.57 & 1.05 \\
GTSRB & benign & 5 & 0.01 & 2 & 1.00 & 0.01 \\
GTSRB & Blend 10\% & 12 & 3.40 & 9 & 0.12 & 0.62 \\
GTSRB & BPP 10\% & 9 & 1.66 & 7 & 0.41 & 0.76 \\
GTSRB & LF 10\% & 12 & 1.68 & 7 & 0.48 & 0.86 \\
GTSRB & TaCT 10\% & 12 & 0.38 & 12 & 0.99 & 0.20 \\
GTSRB & WaNet 10\% & 12 & 1.38 & 8 & 0.29 & 0.59 \\
CIFAR-100 & BadNets 1\% & 12 & 0.79 & 8 & 0.44 & 0.57 \\
CIFAR-100 & BadNets 10\% & 10 & 0.92 & 8 & 0.39 & 0.57 \\
CIFAR-100 & benign & 10 & 0.12 & 6 & 0.99 & 0.06 \\
CIFAR-100 & Blend 1\% & 10 & 1.04 & 7 & 0.23 & 0.62 \\
CIFAR-100 & Blend 10\% & 9 & 1.70 & 5 & 0.16 & 0.71 \\
CIFAR-100 & BPP 10\% & 12 & 1.28 & 8 & 0.20 & 0.75 \\
CIFAR-100 & LF 10\% & 12 & 1.46 & 8 & 0.24 & 0.75 \\
CIFAR-100 & TaCT 10\% & 10 & 0.26 & 9 & 0.98 & 0.13 \\
Tiny ImageNet & BadNets 10\% & 12 & 0.99 & 9 & 0.41 & 0.55 \\
Tiny ImageNet & benign & 9 & 0.04 & 5 & 1.00 & 0.02 \\
Tiny ImageNet & Blend 10\% & 10 & 1.52 & 7 & 0.15 & 0.66 \\
\bottomrule
\end{tabular}
\end{adjustbox}
\end{table}

\FloatBarrier

\section{Comparison with Karayal\c{c}{\i}n et al. in detail}
\label{app:picek}

Karayal\c{c}{\i}n et al.~\cite{karayalcin2026backdoor} remove a backdoor by projecting the direction of 1 selected layer out of every weight matrix that writes to the residual stream, the patch embedding and every attention and MLP output projection. They report attack success falling from \cited{97.7} to \cited{6.7} percent on average. Blend on CIFAR-100 is their 1 survivor. \Cref{tab:erasure} repeats their procedure on our models with their layer selection rule. It adds 2 controls, the same edit restricted to blocks 10 and 11 and the same edit with a random direction of the same norm.

\IfFileExists{tables/erasure.tex}{% generated by scripts/paper/tab_erasure.py from results/_experiments/whole_network_erasure/<folder>.json (10 checkpoints) at 2026-09-11T04:16:29.264278+00:00, commit 4a04e5dd204b14d8368fb7eeef7092bd331e00d9-dirty. Do not edit by hand.
\begin{table}[htbp]
\centering
\small
\caption{Attack success and clean accuracy after removing the backdoor direction from every weight that writes to the residual stream, on ViT-B/16 at 10\% poisoning, or at the rate named in the attack column where it differs.}
\label{tab:erasure}
\begin{adjustbox}{max width=\linewidth}
\begin{tabular}{llrrrrrrr}
\toprule
Dataset & Attack & ASR & CA & Layer & ASR, all writes & CA, all writes & ASR, blocks 10 and 11 & ASR, random direction \\
\midrule
CIFAR-10 & BadNets & 1.000 & 0.929 & 12 & 0.000 & 0.834 & 1.000 & 1.000 \\
CIFAR-10 & Blend & 1.000 & 0.937 & 12 & 0.000 & 0.864 & 1.000 & 1.000 \\
CIFAR-10 & BPP & 0.993 & 0.944 & 12 & 0.000 & 0.883 & 0.990 & 0.993 \\
CIFAR-10 & LF & 0.998 & 0.939 & 9 & 0.034 & 0.943 & 0.997 & 0.998 \\
CIFAR-10 & WaNet & 0.876 & 0.933 & 7 & 0.009 & 0.931 & 0.855 & 0.884 \\
CIFAR-100 & BadNets & 1.000 & 0.807 & 11 & 0.009 & 0.810 & 1.000 & 1.000 \\
CIFAR-100 & Blend (5\%) & 1.000 & 0.833 & 9 & 0.991 & 0.833 & 1.000 & 1.000 \\
CIFAR-100 & Blend & 1.000 & 0.820 & 7 & 1.000 & 0.819 & 1.000 & 1.000 \\
CIFAR-100 & BPP & 0.988 & 0.824 & 12 & 0.000 & 0.820 & 0.985 & 0.988 \\
CIFAR-100 & WaNet & 0.801 & 0.827 & 7 & 0.000 & 0.834 & 0.773 & 0.808 \\
\bottomrule
\end{tabular}
\end{adjustbox}
\end{table}
}{}
\FloatBarrier

The removal reproduces. Over \ErasureCheckpoints{} checkpoints attack success falls from \ErasureAsrBeforeMean{} to \ErasureAsrAfterAllWritesMean{} and clean accuracy from \ErasureCaBeforeMean{} to \ErasureCaAfterAllWritesMean{}. Blend on CIFAR-100 stays at \ErasureAsrAfterAllWritesMax{} at both poison rates exactly as they report. Neither control moves any attack: \ErasureAsrAfterBlocksOneZeroOneOneMean{} for the blocks 10 and 11 edit and \ErasureAsrAfterRandomDirectionMean{} for the random direction. Our first attempt had edited only the last 2 blocks and left every attack intact, which the blocks 10 and 11 column now shows to be expected. The direction is re-created by every block, so removing it means removing it everywhere.

Their second claim, that distributed triggers reach the class token earlier than patch triggers, is measured independently in \cref{fig:tac-layers}. It holds on CIFAR-100 and Tiny ImageNet, where the onset of the backdoor direction comes earlier for Blend, BPP and LF than for BadNets. It reverses on GTSRB, which they did not run.

Their third contribution is a data-free detector that scores the classifier head's rows against the early layers' output projections and flags a model when 1 class stands out by more than 3 standard deviations, and it is reported to flag WaNet and BPP reliably and patch triggers never. We ran the rule with their grid over the number of layers and the floor of the standard deviation on \WeightDetectorBackdoored{} of our backdoored models and \WeightDetectorBenign{} benign ones, and \cref{tab:weight-detector} gives every model. It names the correct target on all \WeightDetectorWanetModels{} WaNet models, with a best $Z$ of \WeightDetectorWanetBestZMin{} to \WeightDetectorWanetBestZMax{}. On BPP it names the target on \WeightDetectorBppNamed{} of \WeightDetectorBppModels{} models, and it also names the target of the \WeightDetectorLfModels{} LF model. It never names the target of a BadNets, Blend or TaCT model. On the benign models \WeightDetectorBenignFlagMin{} to \WeightDetectorBenignFlagMax{} percent of the grid flags, every flag a false alarm. Their WaNet result reproduces and their BPP result reproduces only in part, and \cref{tab:direction-ablation} says why a weight-space rule reads little.

% generated by scripts/paper/app_weight_detector.py from results/_experiments/weight_detector/summary.json at 2026-09-23T19:47:01.967906+00:00, commit 4087ff7b8dddaf7e969f38f703166a0e490aaacf-dirty. Do not edit by hand.
\begin{table}[htbp]
\centering
\small
\caption{The data-free weight detector of Karayalcin et al. over their grid on our models. Named is whether some grid cell flags the model and names its true target, best Z the largest Z among those cells and flag share the share of the grid that flags, which on a benign model is the false-alarm rate.}
\label{tab:weight-detector}
\begin{adjustbox}{max width=\linewidth}
\begin{tabular}{llrlrr}
\toprule
dataset & attack & rate \% & named & best Z & flag share \\
\midrule
CIFAR-100 & BadNets & 10 & no & -- & 0.04 \\
CIFAR-100 & benign & -- & -- & -- & 0.15 \\
CIFAR-100 & Blend & 10 & no & -- & 0.10 \\
CIFAR-100 & BPP & 1 & yes & 9.9 & 0.08 \\
CIFAR-100 & BPP & 5 & no & -- & 0.06 \\
CIFAR-100 & BPP & 10 & no & -- & 0.06 \\
CIFAR-100 & LF & 10 & yes & 10.0 & 0.05 \\
CIFAR-100 & TaCT & 10 & no & -- & 0.05 \\
CIFAR-100 & WaNet & 5 & yes & 139.3 & 0.27 \\
CIFAR-100 & WaNet & 10 & yes & 69.6 & 0.16 \\
CIFAR-10 & BadNets & 10 & no & -- & 0.13 \\
CIFAR-10 & benign & -- & -- & -- & 0.03 \\
CIFAR-10 & BPP & 10 & yes & 71.4 & 0.07 \\
CIFAR-10 & WaNet & 10 & yes & 40.0 & 0.07 \\
Tiny ImageNet & BadNets & 10 & no & -- & 0.13 \\
Tiny ImageNet & benign & -- & -- & -- & 0.11 \\
Tiny ImageNet & BPP & 10 & no & -- & 0.07 \\
Tiny ImageNet & WaNet & 10 & yes & 34.6 & 0.15 \\
\bottomrule
\end{tabular}
\end{adjustbox}
\end{table}
 The backdoor is a rotated direction rather than a coordinate pattern, so a score built from coordinates of the head and the weights sees only its projection.

\section{Seed replicates}
\label{app:seeds}

\Cref{tab:seeds} gives the 3-seed readings behind \cref{sec:robustness:seeds}. Each row is a backdoored model trained 3 times with different seeds. The columns give the smallest attack success rate among the 3 runs, the AUROC of PSBD-TM on each seed with its standard deviation, the same for PSBD-RD and the mean and standard deviation of their paired gain. PSBD-TM's seed standard deviation averages \SeedSdRecommended{} against \SeedSdPublished{} for PSBD-RD, and the paired gain's own standard deviation reaches \SeedSdGainMax{} on its worst row.

% generated by scripts/paper/tab_seeds.py from results/coverage/coverage.json, results/<folder>[_seed_N]/psbd_metrics.json (14 cells), checkpoints/<folder>_seed_N/args.json at 2026-09-23T14:45:07.243909+00:00, commit 039fc6fea9212f8f6d49218ed14b336aa8f6a0d3-dirty. Do not edit by hand.
\begin{table}[htbp]
\centering
\small
\caption{Seed replicates of PSBD-TM and PSBD-RD on the models trained at 3 seeds. Min ASR is the lowest attack success among the 3 runs.}
\label{tab:seeds}
\begin{adjustbox}{max width=\linewidth}
\begin{tabular}{lllrlrlrrr}
\toprule
Dataset & Attack & Rate & Min ASR & PSBD-TM, seeds 0 / 1 / 2 & sd & PSBD-RD, seeds 0 / 1 / 2 & sd & Gain mean & Gain sd \\
\midrule
CIFAR-100 & BadNets & 0.01 & 0.999 & 0.988 / 0.996 / 0.984 & 0.006 & 0.744 / 0.512 / 0.460 & 0.151 & +0.417 & 0.151 \\
CIFAR-100 & BadNets & 0.05 & 1.000 & 0.994 / 0.996 / 0.950 & 0.026 & 0.826 / 0.493 / 0.274 & 0.278 & +0.449 & 0.259 \\
CIFAR-100 & BadNets & 0.1 & 1.000 & 0.998 / 0.998 / 0.995 & 0.002 & 0.497 / 0.659 / 0.565 & 0.081 & +0.423 & 0.081 \\
CIFAR-100 & Blend & 0.01 & 0.988 & 0.976 / 0.991 / 0.923 & 0.036 & 0.792 / 0.991 / 0.871 & 0.100 & +0.079 & 0.095 \\
CIFAR-100 & Blend & 0.05 & 0.999 & 0.979 / 0.995 / 0.982 & 0.009 & 0.999 / 0.996 / 0.902 & 0.055 & +0.020 & 0.053 \\
CIFAR-100 & Blend & 0.1 & 0.999 & 0.998 / 0.996 / 0.989 & 0.005 & 1.000 / 0.995 / 0.992 & 0.004 & -0.001 & 0.003 \\
Tiny ImageNet & BadNets & 0.01 & 1.000 & 0.980 / 0.976 / 0.985 & 0.004 & 0.892 / 0.828 / 0.908 & 0.042 & +0.104 & 0.038 \\
Tiny ImageNet & BadNets & 0.05 & 1.000 & 0.984 / 0.992 / 0.975 & 0.008 & 0.837 / 0.906 / 0.696 & 0.107 & +0.171 & 0.099 \\
Tiny ImageNet & BadNets & 0.1 & 1.000 & 0.993 / 0.986 / 0.969 & 0.012 & 0.973 / 0.856 / 0.650 & 0.164 & +0.156 & 0.151 \\
Tiny ImageNet & Blend & 0.01 & 1.000 & 0.992 / 0.999 / 0.996 & 0.003 & 0.996 / 0.999 / 0.997 & 0.001 & -0.001 & 0.002 \\
Tiny ImageNet & Blend & 0.05 & 0.997 & 0.997 / 0.998 / 0.979 & 0.011 & 0.999 / 1.000 / 0.994 & 0.003 & -0.006 & 0.007 \\
Tiny ImageNet & Blend & 0.1 & 1.000 & 0.998 / 0.997 / 0.997 & 0.001 & 0.999 / 1.000 / 1.000 & 0.000 & -0.002 & 0.001 \\
Tiny ImageNet & WaNet & 0.05 & 0.828 & 0.930 / 0.950 / 0.874 & 0.039 & 0.953 / 0.961 / 0.877 & 0.047 & -0.013 & 0.011 \\
Tiny ImageNet & WaNet & 0.1 & 0.931 & 0.950 / 0.967 / 0.948 & 0.010 & 0.966 / 0.974 / 0.928 & 0.024 & -0.001 & 0.018 \\
\bottomrule
\end{tabular}
\end{adjustbox}
\end{table}

\FloatBarrier

\section{Every detector at both false-positive budgets}
\label{app:detectors}

\Cref{tab:detectors-tpr10,tab:detectors-tpr20} give the 2 false-positive budgets behind \cref{sec:robustness:detectors}, in the same layout as \cref{tab:detectors-auroc}: rows are attacks and columns defenses, ours first, best in each row bold and second best underlined. A threshold is set on the \HeldoutSize{} clean validation images so that the named share of them is flagged, so every column is read at the same budget as every other. Gray marks a reading below chance.

% generated by scripts/paper/tab_detectors.py from results/coverage/coverage.json (69 of 71 clearing cells carry every defense), results/*/detectors/*_metrics.json, results/*/psbd_metrics.json at 2026-09-23T16:58:42.490930+00:00, commit d4a8014354b8fcc8d014beec50fef60ed46fa47d-dirty. Do not edit by hand.
\begin{table*}[htbp]
\centering
\small
\caption{TPR at a 10\% false-positive budget of every defense on every backdoored ViT-B/16 model that carries all 13 of them, 69 of 71 models over 6 datasets. Rows are attacks and columns defenses, ours first. Best in each row is bold and second best underlined. AUROC is one-sided and never flipped, so gray marks a reading below chance, where the detector ordered poisoned and clean inputs the wrong way round. Shaded rows are means. Scale-Up$^\dagger$ is the data-limited variant and IBD-PSC$^\ast$ the calibrated one.}
\label{tab:detectors-tpr10}
\begin{adjustbox}{max width=\textwidth}
\begin{tabular}{lrrrrrrrrrrrrr}
\toprule
attack, rate \% & PSBD-TM & PSBD-RD & Conf & STRIP & Scale-Up & Scale-Up$^\dagger$ & IBD-PSC & IBD-PSC$^\ast$ & TeCo & CD-L & Beatrix & TED & SentiNet \\
\midrule
\multicolumn{14}{l}{\emph{CIFAR-10}} \\
BadNets 1 & 0.959 & \textcolor{black!45}{0.015} & \textcolor{black!45}{0.000} & 0.652 & \textbf{0.997} & \textcolor{black!45}{0.000} & \textcolor{black!45}{0.000} & 0.967 & 0.685 & 0.891 & \underline{0.981} & 0.569 & \textcolor{black!45}{0.201} \\
BadNets 5 & 0.965 & \textcolor{black!45}{0.031} & 0.641 & 0.875 & \textbf{1.000} & \textcolor{black!45}{0.000} & \textcolor{black!45}{0.018} & 0.819 & 0.966 & 0.972 & \underline{0.998} & 0.517 & \textcolor{black!45}{0.119} \\
BadNets 10 & 0.935 & \textcolor{black!45}{0.113} & \textcolor{black!45}{0.486} & \textbf{1.000} & \textcolor{black!45}{0.000} & \textcolor{black!45}{0.000} & 0.947 & \underline{0.992} & 0.882 & \textbf{1.000} & \textbf{1.000} & \textcolor{black!45}{0.291} & \textcolor{black!45}{0.222} \\
Blend 1 & 0.773 & 0.983 & \textcolor{black!45}{0.002} & 0.951 & 0.637 & \textcolor{black!45}{0.000} & \textcolor{black!45}{0.140} & \underline{0.998} & 0.764 & 0.996 & \textbf{1.000} & 0.556 & \textcolor{black!45}{0.000} \\
Blend 5 & 0.908 & 0.973 & 0.633 & 0.885 & 0.553 & \textcolor{black!45}{0.000} & 0.670 & 0.997 & 0.689 & \underline{0.998} & \textbf{1.000} & 0.947 & \textcolor{black!45}{0.000} \\
Blend 10 & 0.754 & \underline{0.980} & \textcolor{black!45}{0.004} & 0.853 & \textcolor{black!45}{0.375} & \textcolor{black!45}{0.000} & \textcolor{black!45}{0.000} & \textbf{1.000} & 0.759 & 0.646 & \textbf{1.000} & 0.917 & \textcolor{black!45}{0.008} \\
BPP 1 & 0.956 & 0.897 & \textcolor{black!45}{0.000} & 0.810 & \textcolor{black!45}{0.010} & \textcolor{black!45}{0.010} & \textcolor{black!45}{0.000} & \underline{0.977} & \textcolor{black!45}{0.011} & 0.877 & \textbf{0.985} & 0.642 & \textcolor{black!45}{0.001} \\
BPP 5 & 0.587 & 0.673 & \textcolor{black!45}{0.000} & 0.760 & \textcolor{black!45}{0.010} & \textcolor{black!45}{0.008} & \textcolor{black!45}{0.000} & \underline{0.955} & \textcolor{black!45}{0.014} & 0.856 & \textbf{0.996} & 0.615 & \textcolor{black!45}{0.000} \\
BPP 10 & 0.736 & 0.720 & \textcolor{black!45}{0.013} & 0.881 & \textcolor{black!45}{0.001} & \textcolor{black!45}{0.001} & \textcolor{black!45}{0.041} & \underline{0.990} & \textcolor{black!45}{0.003} & 0.944 & \textbf{0.999} & 0.761 & \textcolor{black!45}{0.001} \\
LF 1 & \textbf{0.954} & 0.753 & \textcolor{black!45}{0.000} & \textcolor{black!45}{0.368} & \textcolor{black!45}{0.379} & \textcolor{black!45}{0.379} & \textcolor{black!45}{0.000} & 0.673 & \textcolor{black!45}{0.133} & \textcolor{black!45}{0.457} & \underline{0.851} & 0.550 & \textcolor{black!45}{0.000} \\
LF 5 & \textbf{0.983} & 0.837 & \textcolor{black!45}{0.000} & 0.867 & 0.762 & \textcolor{black!45}{0.000} & \textcolor{black!45}{0.000} & 0.875 & \textcolor{black!45}{0.033} & 0.722 & \underline{0.970} & 0.502 & \textcolor{black!45}{0.000} \\
LF 10 & \textbf{0.995} & 0.908 & \textcolor{black!45}{0.125} & 0.762 & 0.642 & \textcolor{black!45}{0.000} & \textcolor{black!45}{0.003} & 0.572 & \textcolor{black!45}{0.086} & 0.596 & \underline{0.941} & 0.643 & \textcolor{black!45}{0.001} \\
SIG 10 & \textcolor{black!45}{0.071} & 0.810 & 0.753 & 0.562 & \textcolor{black!45}{0.437} & \textcolor{black!45}{0.437} & \textcolor{black!45}{0.295} & \textcolor{black!45}{0.138} & \textcolor{black!45}{0.296} & \textcolor{black!45}{0.038} & \underline{0.985} & \textbf{0.995} & \textcolor{black!45}{0.002} \\
TaCT 1 & \textcolor{black!45}{0.005} & \textcolor{black!45}{0.013} & \textcolor{black!45}{0.000} & \textcolor{black!45}{0.089} & \textcolor{black!45}{0.131} & \textcolor{black!45}{0.004} & \textcolor{black!45}{0.000} & \textcolor{black!45}{0.134} & \textcolor{black!45}{0.098} & \textcolor{black!45}{0.005} & \textbf{0.994} & \underline{0.853} & \textcolor{black!45}{0.006} \\
TaCT 5 & \textcolor{black!45}{0.000} & \textcolor{black!45}{0.006} & \textcolor{black!45}{0.098} & \textcolor{black!45}{0.141} & 0.823 & 0.691 & \textcolor{black!45}{0.005} & \textcolor{black!45}{0.270} & \textcolor{black!45}{0.310} & \textcolor{black!45}{0.057} & \textbf{0.987} & \underline{0.950} & \textcolor{black!45}{0.000} \\
TaCT 10 & \textcolor{black!45}{0.030} & \textcolor{black!45}{0.000} & \textcolor{black!45}{0.009} & \textcolor{black!45}{0.373} & \underline{0.950} & \underline{0.950} & \textcolor{black!45}{0.097} & 0.680 & \textcolor{black!45}{0.099} & \textcolor{black!45}{0.499} & \textbf{0.984} & 0.839 & \textcolor{black!45}{0.000} \\
WaNet 5 & 0.838 & \textbf{0.908} & \textcolor{black!45}{0.000} & \textcolor{black!45}{0.089} & \textcolor{black!45}{0.000} & \textcolor{black!45}{0.005} & \textcolor{black!45}{0.000} & \textcolor{black!45}{0.174} & \underline{0.907} & \textcolor{black!45}{0.378} & \textcolor{black!45}{0.334} & 0.559 & \textcolor{black!45}{0.064} \\
WaNet 10 & \textcolor{black!45}{0.087} & \textbf{0.846} & \textcolor{black!45}{0.000} & \textcolor{black!45}{0.115} & \textcolor{black!45}{0.164} & \textcolor{black!45}{0.060} & \textcolor{black!45}{0.001} & \textcolor{black!45}{0.199} & \textcolor{black!45}{0.361} & \textcolor{black!45}{0.304} & 0.533 & \underline{0.656} & \textcolor{black!45}{0.014} \\
\rowcolor{black!8}\emph{CIFAR-10 mean} & 0.641 & 0.581 & \textcolor{black!45}{0.154} & 0.613 & \textcolor{black!45}{0.437} & \textcolor{black!45}{0.141} & \textcolor{black!45}{0.123} & \underline{0.689} & \textcolor{black!45}{0.394} & 0.624 & \textbf{0.919} & 0.687 & \textcolor{black!45}{0.036} \\
\multicolumn{14}{l}{\emph{CIFAR-100}} \\
BadNets 1 & 0.999 & \textcolor{black!45}{0.244} & \textcolor{black!45}{0.017} & \textbf{1.000} & 1.000 & \underline{1.000} & \textcolor{black!45}{0.000} & 0.973 & 0.998 & \textbf{1.000} & \textbf{1.000} & 0.975 & 0.932 \\
BadNets 5 & \underline{0.999} & \textcolor{black!45}{0.335} & 0.958 & \textbf{1.000} & \textbf{1.000} & \textbf{1.000} & 0.841 & 0.984 & 0.998 & \textbf{1.000} & \textcolor{black!45}{0.004} & 0.989 & 0.730 \\
BadNets 10 & \textbf{1.000} & \textcolor{black!45}{0.114} & 0.924 & \underline{0.999} & \textbf{1.000} & \textbf{1.000} & 0.994 & 0.987 & 0.999 & \textbf{1.000} & \textbf{1.000} & 0.941 & 0.895 \\
Blend 1 & 0.959 & \textcolor{black!45}{0.347} & \textcolor{black!45}{0.000} & \textcolor{black!45}{0.274} & \textcolor{black!45}{0.335} & \textcolor{black!45}{0.334} & \textcolor{black!45}{0.126} & 0.979 & \underline{0.986} & \textcolor{black!45}{0.423} & \textbf{0.989} & 0.973 & \textcolor{black!45}{0.003} \\
Blend 5 & 0.941 & 0.999 & 0.960 & 0.736 & \textcolor{black!45}{0.266} & \textcolor{black!45}{0.266} & 0.810 & \underline{1.000} & 0.919 & 0.887 & 0.996 & \textbf{1.000} & \textcolor{black!45}{0.459} \\
Blend 10 & \textbf{1.000} & \underline{1.000} & 0.928 & 0.989 & 0.602 & \textcolor{black!45}{0.348} & 0.999 & \textbf{1.000} & \textcolor{black!45}{0.291} & 0.911 & 0.854 & 0.910 & \textcolor{black!45}{0.008} \\
BPP 1 & 0.849 & \textbf{0.914} & \textcolor{black!45}{0.001} & 0.545 & \textcolor{black!45}{0.264} & \textcolor{black!45}{0.008} & \textcolor{black!45}{0.001} & 0.607 & \textcolor{black!45}{0.076} & \underline{0.901} & \textcolor{black!45}{0.038} & 0.898 & \textcolor{black!45}{0.242} \\
BPP 5 & 0.961 & \underline{0.973} & 0.811 & 0.893 & \textcolor{black!45}{0.002} & \textcolor{black!45}{0.001} & 0.523 & 0.824 & 0.752 & \textbf{0.983} & 0.884 & 0.963 & \textcolor{black!45}{0.029} \\
BPP 10 & \underline{0.963} & 0.960 & \textcolor{black!45}{0.486} & 0.860 & \textcolor{black!45}{0.001} & \textcolor{black!45}{0.001} & \textcolor{black!45}{0.260} & 0.948 & \textcolor{black!45}{0.266} & \textbf{0.977} & 0.923 & 0.960 & \textcolor{black!45}{0.022} \\
LF 1 & \textbf{0.942} & 0.903 & \textcolor{black!45}{0.029} & \textcolor{black!45}{0.477} & \textcolor{black!45}{0.233} & \textcolor{black!45}{0.230} & 0.539 & 0.815 & \textcolor{black!45}{0.247} & \textcolor{black!45}{0.049} & 0.887 & \underline{0.917} & \textcolor{black!45}{0.033} \\
LF 5 & \textbf{0.954} & 0.934 & 0.590 & \textcolor{black!45}{0.299} & \textcolor{black!45}{0.117} & \textcolor{black!45}{0.219} & \textcolor{black!45}{0.031} & \textcolor{black!45}{0.243} & \textcolor{black!45}{0.083} & \textcolor{black!45}{0.030} & \textcolor{black!45}{0.366} & \underline{0.949} & \textcolor{black!45}{0.203} \\
LF 10 & \textbf{0.995} & \underline{0.959} & 0.858 & 0.706 & \textcolor{black!45}{0.163} & \textcolor{black!45}{0.162} & 0.765 & 0.920 & \textcolor{black!45}{0.014} & \textcolor{black!45}{0.010} & 0.688 & 0.950 & \textcolor{black!45}{0.001} \\
TaCT 1 & 0.506 & \textcolor{black!45}{0.023} & \underline{0.632} & \textcolor{black!45}{0.184} & \textcolor{black!45}{0.069} & \textcolor{black!45}{0.184} & 0.575 & \textcolor{black!45}{0.218} & \textcolor{black!45}{0.138} & \textcolor{black!45}{0.000} & \textcolor{black!45}{0.000} & \textbf{0.966} & \textcolor{black!45}{0.046} \\
TaCT 5 & 0.540 & \textcolor{black!45}{0.218} & 0.690 & 0.920 & \textbf{0.977} & \textbf{0.977} & 0.586 & 0.736 & 0.552 & \textcolor{black!45}{0.241} & \textcolor{black!45}{0.000} & \underline{0.943} & \textcolor{black!45}{0.080} \\
TaCT 10 & \textcolor{black!45}{0.161} & \textcolor{black!45}{0.000} & 0.724 & \textcolor{black!45}{0.195} & \underline{0.931} & \underline{0.931} & \textcolor{black!45}{0.345} & 0.690 & \textcolor{black!45}{0.379} & \textcolor{black!45}{0.115} & \textcolor{black!45}{0.000} & \textbf{0.943} & \textcolor{black!45}{0.057} \\
\rowcolor{black!8}\emph{CIFAR-100 mean} & \underline{0.851} & 0.595 & 0.574 & 0.672 & \textcolor{black!45}{0.464} & \textcolor{black!45}{0.444} & \textcolor{black!45}{0.493} & 0.795 & 0.513 & 0.569 & 0.575 & \textbf{0.952} & \textcolor{black!45}{0.249} \\
\multicolumn{14}{l}{\emph{EuroSAT}} \\
Blend 10 & \textcolor{black!45}{0.014} & 0.501 & 0.978 & 0.932 & \underline{1.000} & \textcolor{black!45}{0.000} & 0.768 & 0.578 & \textcolor{black!45}{0.025} & \textcolor{black!45}{0.002} & \textbf{1.000} & 0.947 & \textcolor{black!45}{0.006} \\
SIG 10 & \textcolor{black!45}{0.018} & \textcolor{black!45}{0.094} & \textcolor{black!45}{0.003} & 0.680 & \textcolor{black!45}{0.000} & \underline{0.799} & \textcolor{black!45}{0.120} & 0.774 & \textcolor{black!45}{0.019} & \textcolor{black!45}{0.024} & \textbf{0.992} & \textcolor{black!45}{0.176} & \textcolor{black!45}{0.008} \\
\rowcolor{black!8}\emph{EuroSAT mean} & \textcolor{black!45}{0.016} & \textcolor{black!45}{0.297} & \textcolor{black!45}{0.490} & \underline{0.806} & \textcolor{black!45}{0.500} & \textcolor{black!45}{0.399} & \textcolor{black!45}{0.444} & 0.676 & \textcolor{black!45}{0.022} & \textcolor{black!45}{0.013} & \textbf{0.996} & 0.562 & \textcolor{black!45}{0.007} \\
\multicolumn{14}{l}{\emph{GTSRB}} \\
BadNets 1 & 0.996 & \textcolor{black!45}{0.001} & \textcolor{black!45}{0.000} & \underline{1.000} & \textcolor{black!45}{0.000} & \textbf{1.000} & \textcolor{black!45}{0.000} & 0.702 & 0.952 & 0.940 & \textbf{1.000} & \textbf{1.000} & \textcolor{black!45}{0.014} \\
BadNets 5 & \textbf{1.000} & 0.659 & \textcolor{black!45}{0.000} & \textbf{1.000} & \textcolor{black!45}{0.000} & \textbf{1.000} & \textcolor{black!45}{0.000} & \textbf{1.000} & \underline{0.999} & \textbf{1.000} & \textbf{1.000} & \textbf{1.000} & \textcolor{black!45}{0.035} \\
BadNets 10 & \textbf{1.000} & 0.999 & \textcolor{black!45}{0.000} & \textbf{1.000} & \textcolor{black!45}{0.000} & \textbf{1.000} & \textcolor{black!45}{0.000} & \textbf{1.000} & 0.862 & \underline{1.000} & \textbf{1.000} & \textbf{1.000} & \textcolor{black!45}{0.001} \\
Blend 1 & \textbf{1.000} & \underline{1.000} & \textcolor{black!45}{0.000} & 0.935 & \textcolor{black!45}{0.000} & 0.542 & \textcolor{black!45}{0.000} & \textbf{1.000} & \textcolor{black!45}{0.388} & 0.742 & \textbf{1.000} & \textbf{1.000} & \textcolor{black!45}{0.000} \\
Blend 5 & \underline{1.000} & 0.714 & 0.727 & 0.953 & \textcolor{black!45}{0.000} & 0.867 & 0.912 & 0.994 & 0.750 & 0.961 & \textbf{1.000} & \underline{1.000} & \textcolor{black!45}{0.048} \\
Blend 10 & \underline{1.000} & 0.996 & \textcolor{black!45}{0.000} & 0.936 & \textcolor{black!45}{0.000} & 0.874 & 1.000 & \underline{1.000} & \textcolor{black!45}{0.436} & 1.000 & \textbf{1.000} & \underline{1.000} & \textcolor{black!45}{0.141} \\
BPP 1 & 0.717 & \textcolor{black!45}{0.128} & \textcolor{black!45}{0.000} & \textcolor{black!45}{0.006} & \textcolor{black!45}{0.000} & \textcolor{black!45}{0.005} & \textcolor{black!45}{0.000} & 0.873 & \textcolor{black!45}{0.154} & 0.826 & \textbf{0.964} & \underline{0.895} & \textcolor{black!45}{0.056} \\
BPP 5 & 0.978 & 0.988 & \textcolor{black!45}{0.000} & 0.802 & \textcolor{black!45}{0.000} & \textcolor{black!45}{0.001} & \textcolor{black!45}{0.000} & \textcolor{black!45}{0.438} & \textcolor{black!45}{0.026} & 0.965 & \textbf{0.994} & \underline{0.992} & \textcolor{black!45}{0.067} \\
BPP 10 & 0.971 & 0.971 & \textcolor{black!45}{0.000} & 0.762 & \textcolor{black!45}{0.000} & \textcolor{black!45}{0.000} & \textcolor{black!45}{0.000} & \textcolor{black!45}{0.218} & \textcolor{black!45}{0.066} & 0.919 & \underline{0.995} & \textbf{0.996} & \textcolor{black!45}{0.007} \\
LF 1 & 0.926 & 0.977 & \textcolor{black!45}{0.000} & 0.694 & \textcolor{black!45}{0.000} & 0.830 & \textcolor{black!45}{0.000} & 0.906 & \textcolor{black!45}{0.237} & \textcolor{black!45}{0.047} & \textbf{0.983} & \underline{0.979} & \textcolor{black!45}{0.002} \\
LF 5 & 0.999 & 0.997 & \textcolor{black!45}{0.000} & 0.915 & \textcolor{black!45}{0.000} & 0.781 & \textcolor{black!45}{0.000} & 0.937 & \textcolor{black!45}{0.026} & \textcolor{black!45}{0.075} & \textbf{0.999} & \underline{0.999} & \textcolor{black!45}{0.000} \\
LF 10 & 0.998 & 0.999 & \textcolor{black!45}{0.000} & 0.928 & \textcolor{black!45}{0.000} & 0.750 & \textcolor{black!45}{0.000} & 0.873 & \textcolor{black!45}{0.034} & \textcolor{black!45}{0.076} & \textbf{1.000} & \underline{0.999} & \textcolor{black!45}{0.000} \\
TaCT 5 & 0.502 & \textcolor{black!45}{0.010} & \textcolor{black!45}{0.000} & \textcolor{black!45}{0.426} & \textcolor{black!45}{0.000} & \textbf{1.000} & \textcolor{black!45}{0.094} & 0.724 & \underline{0.993} & \textcolor{black!45}{0.367} & \textbf{1.000} & \textbf{1.000} & \textcolor{black!45}{0.003} \\
TaCT 10 & \textcolor{black!45}{0.213} & \textcolor{black!45}{0.000} & \textcolor{black!45}{0.000} & \textcolor{black!45}{0.207} & \textcolor{black!45}{0.000} & \textbf{1.000} & \textcolor{black!45}{0.000} & 0.919 & \underline{0.924} & \textcolor{black!45}{0.274} & \textcolor{black!45}{0.005} & 0.909 & 0.688 \\
WaNet 10 & 0.949 & \underline{0.950} & \textcolor{black!45}{0.000} & \textcolor{black!45}{0.077} & \textcolor{black!45}{0.000} & \textcolor{black!45}{0.210} & \textcolor{black!45}{0.060} & \textcolor{black!45}{0.220} & 0.925 & \textcolor{black!45}{0.056} & \textcolor{black!45}{0.044} & \textbf{0.953} & \textcolor{black!45}{0.005} \\
\rowcolor{black!8}\emph{GTSRB mean} & \underline{0.883} & 0.692 & \textcolor{black!45}{0.048} & 0.709 & \textcolor{black!45}{0.000} & 0.657 & \textcolor{black!45}{0.138} & 0.787 & 0.518 & 0.617 & 0.866 & \textbf{0.982} & \textcolor{black!45}{0.071} \\
\multicolumn{14}{l}{\emph{SVHN}} \\
Blend 10 & 0.731 & \textcolor{black!45}{0.188} & 0.948 & 0.532 & \textcolor{black!45}{0.097} & \textcolor{black!45}{0.362} & 0.999 & \underline{1.000} & 0.892 & 0.950 & \textbf{1.000} & 0.950 & \textcolor{black!45}{0.000} \\
SIG 5 & \textcolor{black!45}{0.120} & \textcolor{black!45}{0.209} & \textcolor{black!45}{0.003} & \textcolor{black!45}{0.004} & \textcolor{black!45}{0.000} & \textcolor{black!45}{0.009} & \textcolor{black!45}{0.075} & \textcolor{black!45}{0.178} & \textcolor{black!45}{0.455} & \textcolor{black!45}{0.062} & \textbf{0.632} & \textcolor{black!45}{0.287} & \textcolor{black!45}{0.245} \\
\rowcolor{black!8}\emph{SVHN mean} & \textcolor{black!45}{0.425} & \textcolor{black!45}{0.199} & \textcolor{black!45}{0.476} & \textcolor{black!45}{0.268} & \textcolor{black!45}{0.048} & \textcolor{black!45}{0.186} & 0.537 & 0.589 & \underline{0.673} & 0.506 & \textbf{0.816} & 0.618 & \textcolor{black!45}{0.123} \\
\multicolumn{14}{l}{\emph{Tiny ImageNet}} \\
BadNets 1 & 0.989 & 0.700 & \textcolor{black!45}{0.003} & 0.769 & \textbf{0.997} & \textbf{0.997} & 0.758 & 0.916 & \underline{0.992} & \textcolor{black!45}{0.000} & \textcolor{black!45}{0.002} & 0.802 & \textcolor{black!45}{0.192} \\
BadNets 5 & \underline{0.999} & 0.566 & 0.800 & 0.868 & \textbf{0.999} & \underline{0.999} & 0.828 & 0.871 & 0.991 & \textcolor{black!45}{0.000} & \textcolor{black!45}{0.000} & 0.549 & \textcolor{black!45}{0.062} \\
BadNets 10 & \underline{1.000} & 0.951 & 0.960 & 0.769 & \textbf{1.000} & 1.000 & 0.877 & 0.929 & 0.987 & \textcolor{black!45}{0.000} & \textcolor{black!45}{0.000} & \textcolor{black!45}{0.487} & \textcolor{black!45}{0.003} \\
Blend 1 & 0.985 & \underline{0.994} & 0.541 & 0.634 & 0.638 & \textcolor{black!45}{0.273} & 0.966 & \textbf{0.999} & 0.709 & \textcolor{black!45}{0.153} & \textcolor{black!45}{0.001} & \textcolor{black!45}{0.064} & \textcolor{black!45}{0.001} \\
Blend 5 & 0.997 & \underline{0.999} & 0.931 & 0.808 & 0.741 & \textcolor{black!45}{0.388} & 0.998 & \textbf{0.999} & 0.701 & \textcolor{black!45}{0.063} & \textcolor{black!45}{0.000} & \textcolor{black!45}{0.007} & \textcolor{black!45}{0.001} \\
Blend 10 & 0.999 & \underline{0.999} & 0.990 & 0.940 & 0.879 & \textcolor{black!45}{0.000} & 0.999 & \textbf{1.000} & 0.852 & \textcolor{black!45}{0.000} & \textcolor{black!45}{0.000} & \textcolor{black!45}{0.060} & \textcolor{black!45}{0.001} \\
BPP 1 & \textbf{0.972} & \underline{0.972} & \textcolor{black!45}{0.382} & 0.907 & \textcolor{black!45}{0.031} & \textcolor{black!45}{0.029} & 0.959 & 0.970 & \textcolor{black!45}{0.217} & \textcolor{black!45}{0.009} & \textcolor{black!45}{0.002} & \textcolor{black!45}{0.231} & \textcolor{black!45}{0.221} \\
BPP 5 & \textbf{0.984} & 0.977 & 0.779 & 0.873 & \textcolor{black!45}{0.003} & \textcolor{black!45}{0.003} & 0.957 & \underline{0.981} & \textcolor{black!45}{0.100} & \textcolor{black!45}{0.122} & \textcolor{black!45}{0.157} & 0.597 & \textcolor{black!45}{0.312} \\
BPP 10 & \underline{0.981} & \textbf{0.995} & \textcolor{black!45}{0.451} & 0.925 & \textcolor{black!45}{0.007} & \textcolor{black!45}{0.001} & 0.566 & 0.566 & \textcolor{black!45}{0.150} & \textcolor{black!45}{0.053} & \textcolor{black!45}{0.007} & \textcolor{black!45}{0.055} & \textcolor{black!45}{0.122} \\
LF 1 & \textbf{0.929} & \underline{0.918} & \textcolor{black!45}{0.448} & \textcolor{black!45}{0.222} & \textcolor{black!45}{0.430} & \textcolor{black!45}{0.324} & 0.734 & 0.734 & \textcolor{black!45}{0.285} & \textcolor{black!45}{0.283} & \textcolor{black!45}{0.004} & \textcolor{black!45}{0.263} & \textcolor{black!45}{0.011} \\
LF 5 & \textbf{0.987} & \underline{0.987} & 0.893 & 0.617 & 0.565 & \textcolor{black!45}{0.001} & 0.954 & 0.954 & \textcolor{black!45}{0.141} & \textcolor{black!45}{0.011} & \textcolor{black!45}{0.000} & \textcolor{black!45}{0.039} & \textcolor{black!45}{0.003} \\
LF 10 & \textbf{0.973} & \underline{0.970} & 0.838 & \textcolor{black!45}{0.454} & \textcolor{black!45}{0.404} & \textcolor{black!45}{0.404} & 0.910 & 0.910 & \textcolor{black!45}{0.361} & \textcolor{black!45}{0.034} & \textcolor{black!45}{0.002} & \textcolor{black!45}{0.152} & \textcolor{black!45}{0.003} \\
TaCT 1 & \textcolor{black!45}{0.024} & \textcolor{black!45}{0.024} & \textcolor{black!45}{0.190} & \textcolor{black!45}{0.167} & \textcolor{black!45}{0.000} & \textcolor{black!45}{0.000} & \textcolor{black!45}{0.167} & \textcolor{black!45}{0.024} & \textcolor{black!45}{0.000} & \textcolor{black!45}{0.024} & \textcolor{black!45}{0.119} & \textcolor{black!45}{0.143} & \textcolor{black!45}{0.000} \\
TaCT 5 & \textcolor{black!45}{0.381} & \textcolor{black!45}{0.024} & \textcolor{black!45}{0.452} & \textcolor{black!45}{0.071} & \textcolor{black!45}{0.048} & \textcolor{black!45}{0.048} & \textcolor{black!45}{0.310} & \textcolor{black!45}{0.310} & \textcolor{black!45}{0.000} & \textcolor{black!45}{0.000} & \textcolor{black!45}{0.000} & \textcolor{black!45}{0.071} & \textcolor{black!45}{0.000} \\
TaCT 10 & \textcolor{black!45}{0.381} & \textcolor{black!45}{0.048} & \textcolor{black!45}{0.405} & \textcolor{black!45}{0.071} & \textcolor{black!45}{0.024} & \textcolor{black!45}{0.000} & \textcolor{black!45}{0.071} & \textcolor{black!45}{0.071} & \textcolor{black!45}{0.024} & \textcolor{black!45}{0.000} & \textcolor{black!45}{0.000} & \textcolor{black!45}{0.095} & \textcolor{black!45}{0.000} \\
WaNet 5 & \textbf{0.924} & \underline{0.922} & \textcolor{black!45}{0.401} & \textcolor{black!45}{0.105} & \textcolor{black!45}{0.055} & \textcolor{black!45}{0.058} & 0.680 & 0.680 & 0.841 & \textcolor{black!45}{0.070} & \textcolor{black!45}{0.003} & \textcolor{black!45}{0.052} & \textcolor{black!45}{0.011} \\
WaNet 10 & \textbf{0.932} & \underline{0.924} & 0.611 & \textcolor{black!45}{0.112} & \textcolor{black!45}{0.130} & \textcolor{black!45}{0.059} & 0.677 & 0.677 & 0.890 & 0.742 & \textcolor{black!45}{0.001} & \textcolor{black!45}{0.292} & \textcolor{black!45}{0.011} \\
\rowcolor{black!8}\emph{Tiny ImageNet mean} & \textbf{0.849} & \underline{0.763} & 0.593 & 0.548 & \textcolor{black!45}{0.409} & \textcolor{black!45}{0.270} & 0.730 & 0.741 & \textcolor{black!45}{0.485} & \textcolor{black!45}{0.092} & \textcolor{black!45}{0.018} & \textcolor{black!45}{0.233} & \textcolor{black!45}{0.056} \\
\rowcolor{black!8}\emph{all models} & \textbf{0.766} & 0.634 & \textcolor{black!45}{0.349} & 0.626 & \textcolor{black!45}{0.332} & \textcolor{black!45}{0.360} & \textcolor{black!45}{0.378} & \underline{0.743} & \textcolor{black!45}{0.467} & \textcolor{black!45}{0.458} & 0.610 & 0.691 & \textcolor{black!45}{0.097} \\
\bottomrule
\end{tabular}
\end{adjustbox}
\end{table*}

\FloatBarrier
% generated by scripts/paper/tab_detectors.py from results/coverage/coverage.json (69 of 71 clearing cells carry every defense), results/*/detectors/*_metrics.json, results/*/psbd_metrics.json at 2026-09-23T16:58:52.162039+00:00, commit d4a8014354b8fcc8d014beec50fef60ed46fa47d-dirty. Do not edit by hand.
\begin{table*}[htbp]
\centering
\small
\caption{TPR at a 20\% false-positive budget of every defense on every backdoored ViT-B/16 model that carries all 13 of them, 69 of 71 models over 6 datasets. Rows are attacks and columns defenses, ours first. Best in each row is bold and second best underlined. AUROC is one-sided and never flipped, so gray marks a reading below chance, where the detector ordered poisoned and clean inputs the wrong way round. Shaded rows are means. Scale-Up$^\dagger$ is the data-limited variant and IBD-PSC$^\ast$ the calibrated one.}
\label{tab:detectors-tpr20}
\begin{adjustbox}{max width=\textwidth}
\begin{tabular}{lrrrrrrrrrrrrr}
\toprule
attack, rate \% & PSBD-TM & PSBD-RD & Conf & STRIP & Scale-Up & Scale-Up$^\dagger$ & IBD-PSC & IBD-PSC$^\ast$ & TeCo & CD-L & Beatrix & TED & SentiNet \\
\midrule
\multicolumn{14}{l}{\emph{CIFAR-10}} \\
BadNets 1 & 0.978 & \textcolor{black!45}{0.046} & \textcolor{black!45}{0.000} & 0.842 & \underline{0.997} & \underline{0.997} & \textcolor{black!45}{0.000} & 0.985 & 0.879 & 0.914 & \textbf{1.000} & 0.698 & \textcolor{black!45}{0.437} \\
BadNets 5 & 0.995 & \textcolor{black!45}{0.168} & 0.784 & 0.989 & \underline{1.000} & \underline{1.000} & \textcolor{black!45}{0.138} & 0.964 & 0.978 & 0.983 & \textbf{1.000} & 0.672 & \textcolor{black!45}{0.178} \\
BadNets 10 & \underline{1.000} & \textcolor{black!45}{0.301} & 0.792 & \textbf{1.000} & \textbf{1.000} & \textcolor{black!45}{0.000} & 0.998 & \textbf{1.000} & 0.933 & \textbf{1.000} & \textbf{1.000} & \textcolor{black!45}{0.447} & \textcolor{black!45}{0.278} \\
Blend 1 & 0.869 & 0.992 & \textcolor{black!45}{0.109} & 0.979 & 0.767 & 0.523 & 0.934 & \underline{1.000} & 0.835 & 0.998 & \textbf{1.000} & 0.716 & \textcolor{black!45}{0.000} \\
Blend 5 & 0.989 & 0.995 & 0.891 & 0.948 & 0.777 & 0.553 & 0.895 & \underline{1.000} & 0.831 & 0.999 & \textbf{1.000} & 0.979 & \textcolor{black!45}{0.000} \\
Blend 10 & 0.849 & \underline{0.992} & \textcolor{black!45}{0.163} & 0.940 & 0.525 & \textcolor{black!45}{0.375} & \textcolor{black!45}{0.095} & \textbf{1.000} & 0.817 & 0.930 & \textbf{1.000} & 0.978 & \textcolor{black!45}{0.035} \\
BPP 1 & 0.976 & 0.952 & \textcolor{black!45}{0.000} & 0.886 & \textcolor{black!45}{0.035} & \textcolor{black!45}{0.025} & \textcolor{black!45}{0.000} & \underline{0.980} & \textcolor{black!45}{0.060} & 0.916 & \textbf{0.989} & 0.787 & \textcolor{black!45}{0.008} \\
BPP 5 & 0.655 & 0.788 & \textcolor{black!45}{0.002} & 0.854 & \textcolor{black!45}{0.023} & \textcolor{black!45}{0.019} & \textcolor{black!45}{0.001} & \underline{0.978} & \textcolor{black!45}{0.071} & 0.895 & \textbf{0.998} & 0.786 & \textcolor{black!45}{0.002} \\
BPP 10 & 0.830 & 0.799 & \textcolor{black!45}{0.379} & 0.931 & \textcolor{black!45}{0.002} & \textcolor{black!45}{0.002} & \textcolor{black!45}{0.499} & \underline{0.991} & \textcolor{black!45}{0.069} & 0.955 & \textbf{0.999} & 0.838 & \textcolor{black!45}{0.004} \\
LF 1 & \textbf{0.975} & 0.898 & \textcolor{black!45}{0.000} & 0.555 & 0.556 & \textcolor{black!45}{0.462} & \textcolor{black!45}{0.000} & 0.814 & \textcolor{black!45}{0.213} & 0.637 & \underline{0.938} & 0.669 & \textcolor{black!45}{0.005} \\
LF 5 & \textbf{0.994} & 0.947 & \textcolor{black!45}{0.000} & 0.919 & 0.847 & 0.666 & \textcolor{black!45}{0.000} & 0.962 & \textcolor{black!45}{0.057} & 0.855 & \underline{0.992} & 0.638 & \textcolor{black!45}{0.002} \\
LF 10 & \textbf{0.999} & 0.972 & 0.553 & 0.835 & 0.839 & 0.567 & \textcolor{black!45}{0.062} & 0.782 & \textcolor{black!45}{0.125} & 0.838 & \underline{0.983} & 0.756 & \textcolor{black!45}{0.006} \\
SIG 10 & \textcolor{black!45}{0.097} & 0.878 & 0.804 & 0.691 & 0.541 & 0.541 & \textcolor{black!45}{0.327} & \textcolor{black!45}{0.307} & 0.615 & \textcolor{black!45}{0.243} & \underline{0.993} & \textbf{0.997} & \textcolor{black!45}{0.003} \\
TaCT 1 & \textcolor{black!45}{0.009} & \textcolor{black!45}{0.038} & \textcolor{black!45}{0.000} & \textcolor{black!45}{0.207} & \textcolor{black!45}{0.279} & \textcolor{black!45}{0.133} & \textcolor{black!45}{0.000} & \textcolor{black!45}{0.441} & \textcolor{black!45}{0.246} & \textcolor{black!45}{0.064} & \textbf{1.000} & \underline{0.908} & \textcolor{black!45}{0.009} \\
TaCT 5 & \textcolor{black!45}{0.031} & \textcolor{black!45}{0.020} & \textcolor{black!45}{0.317} & \textcolor{black!45}{0.270} & 0.907 & 0.823 & \textcolor{black!45}{0.074} & 0.533 & 0.569 & \textcolor{black!45}{0.285} & \textbf{1.000} & \underline{0.971} & \textcolor{black!45}{0.003} \\
TaCT 10 & \textcolor{black!45}{0.206} & \textcolor{black!45}{0.000} & \textcolor{black!45}{0.152} & 0.575 & \underline{0.992} & \underline{0.992} & \textcolor{black!45}{0.339} & 0.872 & \textcolor{black!45}{0.389} & 0.799 & \textbf{1.000} & 0.908 & \textcolor{black!45}{0.011} \\
WaNet 5 & 0.913 & \textbf{0.965} & \textcolor{black!45}{0.000} & \textcolor{black!45}{0.182} & 0.889 & \textcolor{black!45}{0.008} & \textcolor{black!45}{0.002} & \textcolor{black!45}{0.350} & \underline{0.952} & 0.502 & \textcolor{black!45}{0.491} & 0.689 & \textcolor{black!45}{0.130} \\
WaNet 10 & \textcolor{black!45}{0.268} & \textbf{0.914} & \textcolor{black!45}{0.006} & \textcolor{black!45}{0.215} & \textcolor{black!45}{0.274} & \textcolor{black!45}{0.176} & \textcolor{black!45}{0.027} & \textcolor{black!45}{0.379} & 0.529 & \textcolor{black!45}{0.371} & 0.728 & \underline{0.754} & \textcolor{black!45}{0.036} \\
\rowcolor{black!8}\emph{CIFAR-10 mean} & 0.702 & 0.648 & \textcolor{black!45}{0.275} & 0.712 & 0.625 & \textcolor{black!45}{0.437} & \textcolor{black!45}{0.244} & \underline{0.797} & 0.509 & 0.732 & \textbf{0.951} & 0.788 & \textcolor{black!45}{0.064} \\
\multicolumn{14}{l}{\emph{CIFAR-100}} \\
BadNets 1 & \underline{1.000} & \textcolor{black!45}{0.476} & \textcolor{black!45}{0.305} & \textbf{1.000} & \underline{1.000} & \textbf{1.000} & 0.633 & 1.000 & 1.000 & \textbf{1.000} & \textbf{1.000} & 0.985 & 0.959 \\
BadNets 5 & \textbf{1.000} & 0.672 & 0.986 & \textbf{1.000} & \textbf{1.000} & \textbf{1.000} & 0.933 & 0.997 & \underline{0.999} & \textbf{1.000} & \textcolor{black!45}{0.102} & 0.994 & 0.806 \\
BadNets 10 & \textbf{1.000} & \textcolor{black!45}{0.234} & 0.973 & \textbf{1.000} & \textbf{1.000} & \textbf{1.000} & \textbf{1.000} & \underline{0.999} & \textbf{1.000} & \textbf{1.000} & \textbf{1.000} & 0.979 & 0.936 \\
Blend 1 & 0.983 & 0.598 & \textcolor{black!45}{0.003} & \textcolor{black!45}{0.450} & \textcolor{black!45}{0.335} & \textcolor{black!45}{0.334} & 0.734 & 0.988 & \textbf{0.991} & 0.774 & \underline{0.989} & 0.983 & \textcolor{black!45}{0.010} \\
Blend 5 & 0.976 & \underline{1.000} & 0.995 & 0.861 & \textcolor{black!45}{0.470} & \textcolor{black!45}{0.470} & 0.982 & \underline{1.000} & 0.986 & 0.991 & \underline{1.000} & \textbf{1.000} & 0.884 \\
Blend 10 & \textbf{1.000} & \textbf{1.000} & 0.996 & 0.993 & 0.793 & 0.602 & \underline{1.000} & \textbf{1.000} & 0.845 & 0.999 & 0.991 & 0.996 & \textcolor{black!45}{0.070} \\
BPP 1 & 0.919 & \underline{0.929} & \textcolor{black!45}{0.002} & 0.716 & \textcolor{black!45}{0.276} & \textcolor{black!45}{0.121} & \textcolor{black!45}{0.012} & 0.842 & \textcolor{black!45}{0.160} & \textbf{0.942} & 0.635 & 0.913 & 0.787 \\
BPP 5 & 0.977 & \underline{0.984} & 0.893 & 0.934 & \textcolor{black!45}{0.328} & \textcolor{black!45}{0.004} & 0.865 & 0.957 & 0.851 & \textbf{0.989} & 0.954 & 0.977 & 0.707 \\
BPP 10 & 0.975 & 0.979 & 0.853 & 0.920 & \textcolor{black!45}{0.294} & \textcolor{black!45}{0.003} & 0.864 & 0.984 & \textcolor{black!45}{0.488} & \underline{0.985} & \textbf{0.986} & 0.975 & \textcolor{black!45}{0.064} \\
LF 1 & \textbf{0.947} & 0.940 & \textcolor{black!45}{0.189} & 0.616 & \textcolor{black!45}{0.391} & \textcolor{black!45}{0.388} & 0.728 & 0.917 & \textcolor{black!45}{0.374} & \textcolor{black!45}{0.101} & \underline{0.946} & 0.937 & \textcolor{black!45}{0.101} \\
LF 5 & \textbf{0.959} & 0.955 & 0.693 & \textcolor{black!45}{0.418} & \textcolor{black!45}{0.221} & 0.505 & \textcolor{black!45}{0.146} & 0.521 & \textcolor{black!45}{0.191} & \textcolor{black!45}{0.054} & 0.954 & \underline{0.959} & 0.788 \\
LF 10 & \textbf{0.995} & 0.983 & 0.933 & 0.815 & \textcolor{black!45}{0.163} & \textcolor{black!45}{0.480} & 0.879 & 0.978 & \textcolor{black!45}{0.032} & \textcolor{black!45}{0.033} & \underline{0.995} & 0.976 & \textcolor{black!45}{0.263} \\
TaCT 1 & 0.690 & \textcolor{black!45}{0.103} & 0.736 & \textcolor{black!45}{0.299} & \textcolor{black!45}{0.184} & 0.506 & \underline{0.770} & 0.667 & \textcolor{black!45}{0.287} & \textcolor{black!45}{0.126} & \textcolor{black!45}{0.000} & \textbf{0.977} & \textcolor{black!45}{0.149} \\
TaCT 5 & 0.862 & 0.552 & 0.828 & 0.966 & \textbf{0.989} & \textbf{0.989} & 0.782 & 0.908 & 0.805 & 0.805 & \textcolor{black!45}{0.000} & \underline{0.977} & \textcolor{black!45}{0.241} \\
TaCT 10 & 0.586 & \textcolor{black!45}{0.103} & 0.816 & \textcolor{black!45}{0.425} & \underline{0.943} & \underline{0.943} & 0.747 & 0.931 & 0.644 & 0.552 & \textcolor{black!45}{0.000} & \textbf{0.977} & \textcolor{black!45}{0.230} \\
\rowcolor{black!8}\emph{CIFAR-100 mean} & \underline{0.925} & 0.701 & 0.680 & 0.761 & 0.559 & 0.556 & 0.738 & 0.912 & 0.643 & 0.690 & 0.703 & \textbf{0.974} & \textcolor{black!45}{0.466} \\
\multicolumn{14}{l}{\emph{EuroSAT}} \\
Blend 10 & \textcolor{black!45}{0.251} & 0.739 & 0.984 & 0.964 & \textbf{1.000} & \underline{1.000} & 0.985 & 0.752 & \textcolor{black!45}{0.081} & \textcolor{black!45}{0.003} & \textbf{1.000} & 0.956 & \textcolor{black!45}{0.006} \\
SIG 10 & \textcolor{black!45}{0.135} & \textcolor{black!45}{0.206} & \textcolor{black!45}{0.062} & 0.775 & 0.799 & 0.799 & \textcolor{black!45}{0.237} & \underline{0.858} & \textcolor{black!45}{0.029} & \textcolor{black!45}{0.041} & \textbf{0.999} & \textcolor{black!45}{0.294} & \textcolor{black!45}{0.016} \\
\rowcolor{black!8}\emph{EuroSAT mean} & \textcolor{black!45}{0.193} & \textcolor{black!45}{0.472} & 0.523 & 0.869 & \underline{0.899} & 0.899 & 0.611 & 0.805 & \textcolor{black!45}{0.055} & \textcolor{black!45}{0.022} & \textbf{1.000} & 0.625 & \textcolor{black!45}{0.011} \\
\multicolumn{14}{l}{\emph{GTSRB}} \\
BadNets 1 & \underline{0.999} & \textcolor{black!45}{0.008} & \textcolor{black!45}{0.000} & \textbf{1.000} & \textbf{1.000} & \textbf{1.000} & \textcolor{black!45}{0.003} & 0.936 & 0.988 & 0.992 & \textbf{1.000} & \textbf{1.000} & \textcolor{black!45}{0.014} \\
BadNets 5 & \textbf{1.000} & \underline{0.830} & \textcolor{black!45}{0.075} & \textbf{1.000} & \textbf{1.000} & \textbf{1.000} & \textcolor{black!45}{0.000} & \textbf{1.000} & \textbf{1.000} & \textbf{1.000} & \textbf{1.000} & \textbf{1.000} & \textcolor{black!45}{0.042} \\
BadNets 10 & \textbf{1.000} & \underline{1.000} & \textcolor{black!45}{0.000} & \textbf{1.000} & \textbf{1.000} & \textbf{1.000} & \textcolor{black!45}{0.000} & \textbf{1.000} & 0.960 & \textbf{1.000} & \textbf{1.000} & \textbf{1.000} & \textcolor{black!45}{0.001} \\
Blend 1 & \underline{1.000} & \underline{1.000} & \textcolor{black!45}{0.000} & 0.949 & \textcolor{black!45}{0.000} & 0.718 & \textcolor{black!45}{0.000} & \underline{1.000} & 0.528 & 0.963 & \textbf{1.000} & \underline{1.000} & \textcolor{black!45}{0.000} \\
Blend 5 & \underline{1.000} & 0.834 & 0.844 & 0.966 & \textcolor{black!45}{0.361} & 0.867 & 0.952 & 0.999 & 0.861 & 0.988 & \textbf{1.000} & \underline{1.000} & \textcolor{black!45}{0.064} \\
Blend 10 & \underline{1.000} & 0.998 & \textcolor{black!45}{0.000} & 0.947 & \textcolor{black!45}{0.374} & 0.874 & 1.000 & \underline{1.000} & 0.643 & \textbf{1.000} & \textbf{1.000} & \underline{1.000} & \textcolor{black!45}{0.157} \\
BPP 1 & 0.809 & 0.507 & \textcolor{black!45}{0.000} & \textcolor{black!45}{0.010} & \textcolor{black!45}{0.029} & \textcolor{black!45}{0.017} & \textcolor{black!45}{0.000} & 0.881 & \textcolor{black!45}{0.247} & \underline{0.906} & \textbf{0.983} & 0.903 & \textcolor{black!45}{0.145} \\
BPP 5 & 0.987 & 0.992 & \textcolor{black!45}{0.043} & 0.855 & \textcolor{black!45}{0.001} & \textcolor{black!45}{0.008} & \textcolor{black!45}{0.000} & 0.624 & \textcolor{black!45}{0.067} & 0.981 & \textbf{0.998} & \underline{0.993} & \textcolor{black!45}{0.078} \\
BPP 10 & 0.986 & 0.992 & \textcolor{black!45}{0.002} & 0.824 & \textcolor{black!45}{0.000} & \textcolor{black!45}{0.000} & \textcolor{black!45}{0.000} & \textcolor{black!45}{0.440} & \textcolor{black!45}{0.116} & 0.980 & \textbf{0.998} & \underline{0.996} & \textcolor{black!45}{0.011} \\
LF 1 & 0.973 & \underline{0.985} & \textcolor{black!45}{0.000} & 0.773 & \textcolor{black!45}{0.000} & 0.830 & \textcolor{black!45}{0.001} & 0.928 & \textcolor{black!45}{0.337} & \textcolor{black!45}{0.075} & \textbf{0.988} & 0.982 & \textcolor{black!45}{0.004} \\
LF 5 & \underline{0.999} & \underline{0.999} & \textcolor{black!45}{0.000} & 0.928 & 0.645 & 0.862 & \textcolor{black!45}{0.000} & 0.971 & \textcolor{black!45}{0.038} & \textcolor{black!45}{0.263} & \textbf{1.000} & 0.999 & \textcolor{black!45}{0.000} \\
LF 10 & 0.999 & \underline{0.999} & \textcolor{black!45}{0.000} & 0.939 & \textcolor{black!45}{0.474} & 0.750 & \textcolor{black!45}{0.072} & 0.921 & \textcolor{black!45}{0.052} & \textcolor{black!45}{0.177} & \textbf{1.000} & \textbf{1.000} & \textcolor{black!45}{0.000} \\
TaCT 5 & 0.818 & \textcolor{black!45}{0.041} & \textcolor{black!45}{0.221} & 0.603 & \textcolor{black!45}{0.000} & \textbf{1.000} & \textcolor{black!45}{0.468} & \underline{0.904} & \textbf{1.000} & 0.719 & \textbf{1.000} & \textbf{1.000} & \textcolor{black!45}{0.003} \\
TaCT 10 & 0.843 & \textcolor{black!45}{0.000} & \textcolor{black!45}{0.000} & \textcolor{black!45}{0.369} & \textcolor{black!45}{0.000} & \textbf{1.000} & \textcolor{black!45}{0.000} & 0.960 & \underline{0.988} & 0.638 & 0.640 & 0.977 & 0.795 \\
WaNet 10 & 0.952 & \textbf{0.960} & \textcolor{black!45}{0.007} & \textcolor{black!45}{0.183} & \textcolor{black!45}{0.097} & \textcolor{black!45}{0.211} & \textcolor{black!45}{0.134} & \textcolor{black!45}{0.278} & 0.938 & \textcolor{black!45}{0.168} & \textcolor{black!45}{0.261} & \underline{0.957} & \textcolor{black!45}{0.019} \\
\rowcolor{black!8}\emph{GTSRB mean} & \underline{0.958} & 0.743 & \textcolor{black!45}{0.079} & 0.756 & \textcolor{black!45}{0.332} & 0.676 & \textcolor{black!45}{0.175} & 0.856 & 0.584 & 0.723 & 0.924 & \textbf{0.987} & \textcolor{black!45}{0.089} \\
\multicolumn{14}{l}{\emph{SVHN}} \\
Blend 10 & 0.890 & \textcolor{black!45}{0.358} & 0.964 & 0.650 & \textcolor{black!45}{0.225} & \textcolor{black!45}{0.362} & 0.999 & \underline{1.000} & 0.962 & 0.995 & \textbf{1.000} & 0.985 & \textcolor{black!45}{0.000} \\
SIG 5 & \textcolor{black!45}{0.482} & \textcolor{black!45}{0.339} & \textcolor{black!45}{0.021} & \textcolor{black!45}{0.021} & 0.504 & \textcolor{black!45}{0.320} & \textcolor{black!45}{0.210} & \textcolor{black!45}{0.402} & 0.674 & \textcolor{black!45}{0.162} & \underline{0.852} & \textcolor{black!45}{0.467} & \textbf{0.880} \\
\rowcolor{black!8}\emph{SVHN mean} & 0.686 & \textcolor{black!45}{0.349} & \textcolor{black!45}{0.493} & \textcolor{black!45}{0.335} & \textcolor{black!45}{0.364} & \textcolor{black!45}{0.341} & 0.605 & 0.701 & \underline{0.818} & 0.578 & \textbf{0.926} & 0.726 & \textcolor{black!45}{0.440} \\
\multicolumn{14}{l}{\emph{Tiny ImageNet}} \\
BadNets 1 & \underline{0.997} & 0.836 & \textcolor{black!45}{0.243} & 0.937 & \underline{0.997} & \underline{0.997} & 0.980 & 0.979 & \textbf{0.998} & \textcolor{black!45}{0.000} & \textcolor{black!45}{0.123} & 0.886 & \textcolor{black!45}{0.251} \\
BadNets 5 & \textbf{1.000} & 0.739 & 0.938 & 0.967 & \underline{0.999} & \underline{0.999} & 0.954 & 0.981 & 0.996 & 0.970 & \textcolor{black!45}{0.003} & 0.754 & \textcolor{black!45}{0.095} \\
BadNets 10 & \underline{1.000} & 0.985 & 0.988 & 0.927 & \textbf{1.000} & 1.000 & 0.966 & 0.982 & 0.994 & 0.992 & \textcolor{black!45}{0.000} & 0.665 & \textcolor{black!45}{0.013} \\
Blend 1 & 0.992 & \underline{0.998} & 0.869 & 0.728 & 0.638 & \textcolor{black!45}{0.468} & \textbf{0.999} & \textbf{0.999} & 0.943 & \textcolor{black!45}{0.249} & \textcolor{black!45}{0.241} & \textcolor{black!45}{0.282} & \textcolor{black!45}{0.001} \\
Blend 5 & \underline{0.999} & \textbf{0.999} & 0.983 & 0.867 & 0.741 & \textcolor{black!45}{0.480} & 0.999 & \underline{0.999} & 0.942 & \textcolor{black!45}{0.197} & \textcolor{black!45}{0.000} & \textcolor{black!45}{0.143} & \textcolor{black!45}{0.002} \\
Blend 10 & \underline{0.999} & \textbf{1.000} & 0.993 & 0.963 & 0.879 & 0.660 & \underline{0.999} & \textbf{1.000} & 0.972 & \textcolor{black!45}{0.001} & \textcolor{black!45}{0.000} & \textcolor{black!45}{0.267} & \textcolor{black!45}{0.001} \\
BPP 1 & \textbf{0.975} & \underline{0.974} & 0.781 & 0.936 & \textcolor{black!45}{0.104} & \textcolor{black!45}{0.105} & 0.969 & 0.973 & \textcolor{black!45}{0.454} & \textcolor{black!45}{0.124} & \textcolor{black!45}{0.003} & \textcolor{black!45}{0.422} & \textcolor{black!45}{0.399} \\
BPP 5 & \underline{0.986} & \textbf{0.987} & 0.872 & 0.912 & \textcolor{black!45}{0.009} & \textcolor{black!45}{0.010} & 0.966 & 0.986 & \textcolor{black!45}{0.233} & 0.650 & \textcolor{black!45}{0.190} & 0.776 & \textcolor{black!45}{0.381} \\
BPP 10 & \underline{0.991} & \textbf{0.996} & 0.815 & 0.952 & \textcolor{black!45}{0.020} & \textcolor{black!45}{0.007} & 0.832 & 0.832 & \textcolor{black!45}{0.341} & 0.900 & \textcolor{black!45}{0.248} & \textcolor{black!45}{0.221} & \textcolor{black!45}{0.187} \\
LF 1 & \textbf{0.935} & \underline{0.934} & 0.610 & \textcolor{black!45}{0.360} & \textcolor{black!45}{0.430} & \textcolor{black!45}{0.430} & 0.846 & 0.846 & \textcolor{black!45}{0.418} & 0.595 & \textcolor{black!45}{0.008} & \textcolor{black!45}{0.434} & \textcolor{black!45}{0.029} \\
LF 5 & \underline{0.988} & \textbf{0.989} & 0.925 & 0.748 & 0.565 & \textcolor{black!45}{0.311} & 0.967 & 0.967 & \textcolor{black!45}{0.222} & \textcolor{black!45}{0.048} & \textcolor{black!45}{0.001} & \textcolor{black!45}{0.157} & \textcolor{black!45}{0.008} \\
LF 10 & \textbf{0.977} & \underline{0.976} & 0.877 & 0.589 & \textcolor{black!45}{0.404} & 0.643 & 0.943 & 0.943 & \textcolor{black!45}{0.483} & \textcolor{black!45}{0.209} & \textcolor{black!45}{0.054} & \textcolor{black!45}{0.339} & \textcolor{black!45}{0.005} \\
TaCT 1 & \textcolor{black!45}{0.119} & \textcolor{black!45}{0.071} & \textcolor{black!45}{0.405} & \textcolor{black!45}{0.310} & \textcolor{black!45}{0.000} & \textcolor{black!45}{0.000} & \textcolor{black!45}{0.333} & \textcolor{black!45}{0.167} & \textcolor{black!45}{0.024} & \textcolor{black!45}{0.119} & \textcolor{black!45}{0.143} & \textbf{0.571} & \textcolor{black!45}{0.071} \\
TaCT 5 & \underline{0.500} & \textcolor{black!45}{0.119} & \textbf{0.595} & \textcolor{black!45}{0.262} & \textcolor{black!45}{0.048} & \textcolor{black!45}{0.048} & \underline{0.500} & \underline{0.500} & \textcolor{black!45}{0.071} & \textcolor{black!45}{0.000} & \textcolor{black!45}{0.000} & \textcolor{black!45}{0.357} & \textcolor{black!45}{0.024} \\
TaCT 10 & \textbf{0.619} & \textcolor{black!45}{0.071} & \underline{0.500} & \textcolor{black!45}{0.167} & \textcolor{black!45}{0.024} & \textcolor{black!45}{0.024} & \textcolor{black!45}{0.190} & \textcolor{black!45}{0.190} & \textcolor{black!45}{0.071} & \textcolor{black!45}{0.095} & \textcolor{black!45}{0.000} & \textcolor{black!45}{0.357} & \textcolor{black!45}{0.024} \\
WaNet 5 & \underline{0.928} & \textbf{0.938} & 0.521 & \textcolor{black!45}{0.218} & \textcolor{black!45}{0.117} & \textcolor{black!45}{0.125} & 0.771 & 0.771 & 0.889 & \textcolor{black!45}{0.394} & \textcolor{black!45}{0.006} & \textcolor{black!45}{0.226} & \textcolor{black!45}{0.037} \\
WaNet 10 & \textbf{0.970} & \underline{0.963} & 0.727 & \textcolor{black!45}{0.228} & \textcolor{black!45}{0.130} & \textcolor{black!45}{0.132} & 0.754 & 0.754 & 0.933 & 0.945 & \textcolor{black!45}{0.002} & 0.552 & \textcolor{black!45}{0.032} \\
\rowcolor{black!8}\emph{Tiny ImageNet mean} & \textbf{0.881} & 0.799 & 0.744 & 0.651 & \textcolor{black!45}{0.418} & \textcolor{black!45}{0.379} & \underline{0.822} & 0.816 & 0.587 & \textcolor{black!45}{0.382} & \textcolor{black!45}{0.060} & \textcolor{black!45}{0.436} & \textcolor{black!45}{0.092} \\
\rowcolor{black!8}\emph{all models} & \underline{0.835} & 0.703 & \textcolor{black!45}{0.450} & 0.711 & \textcolor{black!45}{0.496} & 0.511 & \textcolor{black!45}{0.500} & \textbf{0.837} & 0.570 & 0.610 & 0.673 & 0.778 & \textcolor{black!45}{0.173} \\
\bottomrule
\end{tabular}
\end{adjustbox}
\end{table*}

\FloatBarrier

\section{Detector acceptance smoke test}
\label{app:smoke}

\Cref{tab:detector-smoke} is the acceptance test each port had to pass before it entered the comparison above: a benign, a BadNets and a Blend GTSRB model at 5\% poisoning with 500 validation images, on the same splits and threshold rule as PSBD. The passes column is the forward passes each detector spends per input and the last column the measured seconds per input on 1 A100. Every detector reads chance on the benign model, which shows it has no signal where there is no backdoor. It cannot check the sign convention, since a benign model reads chance under either sign, and the sign is checked instead by the backdoored models of the comparison. The numbers here are not the comparison, which \cref{tab:detectors-auroc} carries over every backdoored model.

% generated by scripts/paper/tab_detector_smoke.py from docs/runs/2026-09-10-detector-smoke.md at 2026-09-10T23:37:23.095504+00:00, commit 831df7ef5e8bc9285264b261f848e1f2cad5ef61-dirty. Do not edit by hand.
\begin{table}[t]
\centering
\small
\caption{Competitor detectors on the login-node smoke run: 3 GTSRB checkpoints at the middle panel rate, AUROC at the headline quantile on the benign reference and the 2 backdoored models, TPR at the smallest budget on the backdoored ones, forward passes per input and seconds per input. Acceptance only, the panel sweep is pending.}
\label{tab:detector-smoke}
\begin{tabular}{lrrrrrrr}
\toprule
detector & passes & AUROC benign & AUROC BadNet & AUROC Blend & TPR@1\% BadNet & TPR@1\% Blend & s per input \\
\midrule
confidence & 1 & 0.503 & 0.738 & 0.876 & 0.000 & 0.034 & 0.001 \\
strip & 8 & 0.500 & 1.000 & 0.976 & 1.000 & 0.905 & 0.005 \\
scale\_up & 6 & 0.499 & 0.914 & 0.689 & 0.000 & 0.000 & 0.004 \\
scale\_up\_data\_limited & 6 & 0.499 & 0.909 & 0.678 & 0.000 & 0.000 & 0.002 \\
ibd\_psc & 6 & 0.500 & 0.481 & 0.970 & 0.000 & 0.724 & 0.002 \\
teco & 71 & 0.495 & 0.994 & 0.922 & 0.757 & 0.258 & 0.057 \\
cd\_l & 251 & 0.493 & 1.000 & 0.962 & 1.000 & 0.596 & 0.112 \\
beatrix & 1 & 0.510 & 0.453 & 0.819 & 0.000 & 0.000 & 0.001 \\
ted & 1 & 0.509 & 0.999 & 0.990 & 0.990 & 0.811 & 0.001 \\
sentinet & 202 & 0.499 & 0.107 & 0.379 & 0.056 & 0.051 & 0.057 \\
\bottomrule
\end{tabular}
\end{table}

\FloatBarrier

\section{The adaptive attacker per model}
\label{app:adaptive}

\Cref{tab:adaptive-attacker} gives every evasive model behind \cref{tab:adaptive}: its attack success after the evasion, its clean-accuracy cost, the AUROC of the probe the attacker trained against before and after, the mean AUROC of the 2 probes it never saw, and the AUROC of the min-rank union of all 3.

% generated by scripts/paper/tab_adaptive.py from results/adaptive_attacker_analysis.json, results/multi_probe_analysis.json at 2026-09-11T04:21:55.153420+00:00, commit ce6f85eb8836cbaebfd6366152eb1269e7de2cb9-dirty. Do not edit by hand.
\begin{table}[htbp]
\centering
\small
\caption{The adaptive attacker on the models whose evasive version keeps attack success at or above 0.9.}
\label{tab:adaptive-attacker}
\begin{adjustbox}{max width=\linewidth}
\begin{tabular}{llllrrrrrr}
\toprule
arch & dataset & attack & rate & evade ASR & dCA & probed base & probed evade & transfer mean & union \\
\midrule
vit & CIFAR-100 & Adaptive-Blend & 0.05 & 0.948 & -0.026 & 0.927 & 0.034 & 0.936 & 0.944 \\
vit & CIFAR-100 & Adaptive-Blend & 0.1 & 0.925 & -0.045 & 0.951 & 0.039 & 0.921 & 0.927 \\
vit & CIFAR-100 & BadNets & 0.01 & 0.999 & -0.018 & 0.960 & 0.048 & 0.806 & 0.984 \\
vit & CIFAR-100 & BadNets & 0.05 & 1.000 & -0.026 & 0.977 & 0.004 & 0.785 & 0.994 \\
vit & CIFAR-100 & BadNets & 0.1 & 1.000 & -0.018 & 0.989 & 0.005 & 0.846 & 0.985 \\
vit & CIFAR-100 & Blend & 0.01 & 0.999 & -0.024 & 0.945 & 0.680 & 0.961 & 0.998 \\
vit & CIFAR-100 & Blend & 0.05 & 1.000 & -0.032 & 0.989 & 0.807 & 0.978 & 0.983 \\
vit & CIFAR-100 & Blend & 0.1 & 1.000 & -0.029 & 0.982 & 0.012 & 0.798 & 0.991 \\
vit & CIFAR-100 & BPP & 0.01 & 0.961 & -0.023 & 0.927 & 0.671 & 0.897 & 0.968 \\
vit & CIFAR-100 & BPP & 0.05 & 0.977 & -0.035 & 0.981 & 0.275 & 0.952 & 0.982 \\
vit & CIFAR-100 & BPP & 0.1 & 0.996 & -0.045 & 0.993 & 0.204 & 0.909 & 0.981 \\
vit & CIFAR-100 & LF & 0.01 & 0.938 & -0.042 & 0.944 & 0.410 & 0.838 & 0.852 \\
vit & CIFAR-100 & LF & 0.05 & 0.989 & -0.032 & 0.957 & 0.060 & 0.936 & 0.986 \\
vit & CIFAR-100 & LF & 0.1 & 0.996 & -0.026 & 0.994 & 0.014 & 0.938 & 0.967 \\
\bottomrule
\end{tabular}
\end{adjustbox}
\end{table}

\FloatBarrier

\section{Swin-S per model}
\label{app:swin}

\Cref{tab:swin-attacks} gives every Swin-S model behind \cref{sec:swin} separately, and \cref{tab:swin-placements} every placement we cached on Swin with the number of models behind each. \Cref{tab:swin-gains} gives the paired comparisons with their bootstrap intervals.

% generated by scripts/paper/tab_swin.py from results/swin_*/psbd_metrics.json (90 implanted cells), checkpoints/swin_*/args.json at 2026-09-23T21:39:54.462617+00:00, commit 33809450714a8a0518a06b5ca5e2d64dacb27d3a-dirty. Do not edit by hand.
\begin{table}[htbp]
\centering
\small
\caption{Swin-S per attack. Token masking at the attention input against dropout after both residual adds, at the adaptive rule and the headline quantile. The shaded row is the mean over the dataset above it.}
\label{tab:swin-attacks}
\begin{adjustbox}{max width=\linewidth}
\begin{tabular}{llrrrr}
\toprule
attack & rate \% & AUROC & TPR@10 & TPR@20 & AUROC published \\
\midrule
\multicolumn{6}{l}{\emph{CIFAR-10}} \\
Adaptive-Blend & 5 & 0.496 & 0.03 & 0.10 & 0.787 \\
Adaptive-Blend & 10 & 0.992 & 0.99 & 0.99 & 0.981 \\
BadNets & 0.5 & 0.998 & 0.99 & 1.00 & 0.942 \\
BadNets & 1 & 0.993 & 1.00 & 1.00 & 0.839 \\
BadNets & 5 & 1.000 & 1.00 & 1.00 & 0.500 \\
BadNets & 10 & 1.000 & 1.00 & 1.00 & 0.996 \\
Blend & 0.5 & 0.991 & 0.97 & 1.00 & 1.000 \\
Blend & 1 & 0.990 & 0.97 & 0.99 & 0.999 \\
Blend & 5 & 0.997 & 1.00 & 1.00 & 0.994 \\
Blend & 10 & 1.000 & 1.00 & 1.00 & 1.000 \\
BPP & 0.5 & 0.985 & 0.96 & 0.98 & 0.989 \\
BPP & 1 & 0.800 & 0.38 & 0.60 & 0.995 \\
BPP & 5 & 0.997 & 0.99 & 0.99 & 0.999 \\
BPP & 10 & 0.998 & 1.00 & 1.00 & 0.998 \\
Label-Consistent & 10 & 0.965 & 0.95 & 0.96 & 0.530 \\
LF & 0.5 & 0.859 & 0.32 & 0.81 & 0.693 \\
LF & 1 & 0.996 & 0.99 & 0.99 & 0.989 \\
LF & 5 & 0.940 & 0.80 & 0.88 & 0.813 \\
LF & 10 & 0.982 & 0.93 & 0.96 & 0.953 \\
SIG & 5 & -- & -- & -- & 0.977 \\
SIG & 10 & 0.954 & 0.93 & 0.94 & 0.937 \\
TaCT & 1 & -- & -- & -- & -- \\
TaCT & 5 & -- & -- & -- & 0.999 \\
TaCT & 10 & -- & -- & -- & -- \\
WaNet & 5 & 0.980 & 0.85 & 0.98 & 0.961 \\
WaNet & 10 & 0.952 & 0.69 & 0.96 & 0.927 \\
\rowcolor{black!8}\emph{mean} &  & 0.948 & 0.85 & 0.92 & 0.908 \\
\multicolumn{6}{l}{\emph{CIFAR-100}} \\
Adaptive-Blend & 5 & 0.921 & 0.91 & 0.91 & 0.920 \\
Adaptive-Blend & 10 & 0.970 & 0.97 & 0.97 & 0.973 \\
BadNets & 0.5 & 1.000 & 1.00 & 1.00 & 0.917 \\
BadNets & 1 & 1.000 & 1.00 & 1.00 & 0.643 \\
BadNets & 5 & 1.000 & 1.00 & 1.00 & 0.853 \\
BadNets & 10 & 1.000 & 1.00 & 1.00 & 0.927 \\
Blend & 0.5 & 0.986 & 0.96 & 0.99 & 0.993 \\
Blend & 1 & 0.999 & 1.00 & 1.00 & 0.992 \\
Blend & 5 & 1.000 & 1.00 & 1.00 & 0.998 \\
Blend & 10 & 0.999 & 1.00 & 1.00 & 0.999 \\
BPP & 0.5 & 0.985 & 0.96 & 0.98 & 0.977 \\
BPP & 1 & 0.997 & 0.99 & 1.00 & 0.994 \\
BPP & 5 & 0.998 & 1.00 & 1.00 & 0.995 \\
BPP & 10 & 0.998 & 0.99 & 1.00 & 0.996 \\
Label-Consistent & 1 & 0.999 & 1.00 & 1.00 & 0.716 \\
LF & 0.5 & 0.959 & 0.90 & 0.96 & 0.966 \\
LF & 1 & 0.945 & 0.85 & 0.95 & 0.874 \\
LF & 5 & 0.991 & 0.99 & 0.99 & 0.957 \\
LF & 10 & 0.996 & 0.99 & 1.00 & 0.986 \\
TaCT & 1 & 0.958 & 0.84 & 0.97 & 0.668 \\
TaCT & 5 & 0.805 & 0.21 & 0.44 & 0.597 \\
TaCT & 10 & 0.960 & 0.95 & 0.98 & 0.516 \\
WaNet & 5 & 0.904 & 0.87 & 0.88 & 0.202 \\
WaNet & 10 & 0.927 & 0.89 & 0.93 & 0.803 \\
\rowcolor{black!8}\emph{mean} &  & 0.971 & 0.93 & 0.95 & 0.853 \\
\multicolumn{6}{l}{\emph{GTSRB}} \\
Adaptive-Blend & 5 & 0.925 & 0.88 & 0.89 & 0.926 \\
Adaptive-Blend & 10 & 1.000 & 1.00 & 1.00 & 0.995 \\
BadNets & 0.5 & 1.000 & 1.00 & 1.00 & 0.200 \\
BadNets & 1 & 1.000 & 1.00 & 1.00 & 0.729 \\
BadNets & 5 & 1.000 & 1.00 & 1.00 & 0.588 \\
BadNets & 10 & 1.000 & 1.00 & 1.00 & 0.989 \\
Blend & 0.5 & 0.987 & 0.96 & 0.97 & 1.000 \\
Blend & 1 & 0.978 & 0.92 & 0.94 & 0.999 \\
Blend & 5 & 1.000 & 1.00 & 1.00 & 1.000 \\
Blend & 10 & 0.987 & 0.95 & 0.97 & 1.000 \\
BPP & 0.5 & 0.642 & 0.36 & 0.43 & 0.703 \\
BPP & 1 & 0.979 & 0.93 & 0.99 & 0.991 \\
BPP & 5 & 0.997 & 0.99 & 1.00 & 1.000 \\
BPP & 10 & 0.998 & 0.99 & 1.00 & 0.999 \\
Label-Consistent & 5 & 0.996 & 0.99 & 0.99 & 0.486 \\
Label-Consistent & 10 & 0.998 & 1.00 & 1.00 & 0.384 \\
LF & 0.5 & 0.975 & 0.95 & 0.97 & 0.711 \\
LF & 1 & 0.965 & 0.89 & 0.98 & 0.587 \\
LF & 5 & 0.980 & 0.94 & 0.97 & 0.670 \\
LF & 10 & 0.998 & 1.00 & 1.00 & 0.797 \\
WaNet & 5 & 0.885 & 0.87 & 0.87 & 0.301 \\
WaNet & 10 & 0.919 & 0.83 & 0.96 & 0.458 \\
\rowcolor{black!8}\emph{mean} &  & 0.964 & 0.93 & 0.95 & 0.751 \\
\multicolumn{6}{l}{\emph{Tiny ImageNet}} \\
Adaptive-Blend & 10 & 0.987 & 0.99 & 0.99 & 0.987 \\
BadNets & 1 & 0.998 & 1.00 & 1.00 & 0.651 \\
BadNets & 5 & 0.999 & 1.00 & 1.00 & 0.690 \\
BadNets & 10 & 1.000 & 1.00 & 1.00 & 0.848 \\
Blend & 1 & 1.000 & 1.00 & 1.00 & 0.998 \\
Blend & 5 & 1.000 & 1.00 & 1.00 & 0.998 \\
Blend & 10 & 1.000 & 1.00 & 1.00 & 0.999 \\
BPP & 1 & 0.997 & 1.00 & 1.00 & 0.996 \\
BPP & 5 & 0.999 & 1.00 & 1.00 & 0.999 \\
BPP & 10 & 1.000 & 1.00 & 1.00 & 0.999 \\
LF & 1 & 0.947 & 0.93 & 0.93 & 0.949 \\
LF & 5 & 0.995 & 1.00 & 1.00 & 0.993 \\
LF & 10 & 0.998 & 1.00 & 1.00 & 0.994 \\
TaCT & 1 & 0.517 & 0.40 & 0.60 & -- \\
TaCT & 5 & 0.520 & 0.48 & 0.60 & -- \\
TaCT & 10 & 0.504 & 0.14 & 0.29 & -- \\
WaNet & 5 & 0.991 & 0.99 & 0.99 & 0.834 \\
WaNet & 10 & 0.992 & 0.99 & 0.99 & 0.431 \\
\rowcolor{black!8}\emph{mean} &  & 0.914 & 0.88 & 0.91 & 0.891 \\
\bottomrule
\end{tabular}
\end{adjustbox}
\end{table}

\FloatBarrier
% generated by scripts/paper/tab_swin.py from results/swin_*/psbd_metrics.json (83 implanted cells), checkpoints/swin_*/args.json at 2026-09-23T11:34:10.367602+00:00, commit 2fb3169f973b846e39a6f585b342a893fa1af109-dirty. Do not edit by hand.
\begin{table}[htbp]
\centering
\small
\caption{Swin-S: mean AUROC at the headline quantile for every cached placement with at least 5 implanted models, at the matched and adaptive rules. Models are Swin folders whose sidecar attack success clears the bar, 25 on CIFAR-10, 21 on CIFAR-100, 22 on GTSRB, 15 on Tiny ImageNet. Coverage is uneven across placements, so rows are not paired.}
\label{tab:swin-placements}
\begin{adjustbox}{max width=\linewidth}
\begin{tabular}{lrrrrr}
\toprule
placement & n match & AUROC matched & n adapt & AUROC adaptive & below chance \\
\midrule
after\_embedding & 82 & 0.821 & 82 & 0.852 & 9 \\
before\_attention & 82 & 0.872 & 82 & 0.890 & 6 \\
post\_residual & 82 & 0.779 & 82 & 0.845 & 10 \\
pre\_residual\_blocks\_1\_8 & 82 & 0.770 & 82 & 0.845 & 10 \\
before\_attention\_norm & 81 & 0.810 & 81 & 0.870 & 6 \\
pre\_residual & 81 & 0.802 & 81 & 0.842 & 10 \\
after\_attention\_residual & 80 & 0.786 & 80 & 0.923 & 1 \\
after\_attention\_residual\_token\_mask & 80 & 0.835 & 80 & 0.795 & 12 \\
before\_attention\_norm\_channel\_mask & 80 & 0.832 & 80 & 0.881 & 5 \\
before\_attention\_norm\_gaussian & 80 & 0.902 & 80 & 0.927 & 3 \\
before\_attention\_norm\_token\_mask & 80 & 0.956 & 80 & 0.969 & 1 \\
before\_attention\_residual\_token\_mask & 80 & 0.799 & 80 & 0.834 & 10 \\
before\_mlp\_gaussian & 80 & 0.788 & 80 & 0.725 & 21 \\
before\_mlp\_norm & 80 & 0.783 & 80 & 0.787 & 13 \\
before\_mlp\_norm\_token\_mask & 80 & 0.709 & 80 & 0.808 & 11 \\
both\_sublayer\_inputs\_token\_mask & 80 & 0.894 & 80 & 0.946 & 1 \\
input\_pixels\_scale\_up & 80 & 0.764 & 80 & 0.813 & 10 \\
mlp\_norm\_out\_gain\_scale & 80 & 0.889 & 80 & 0.814 & 12 \\
pre\_residual\_blocks\_9\_16 & 82 & 0.839 & 80 & 0.882 & 4 \\
before\_attention\_norm\_blocks\_1\_8\_token\_mask & 80 & 0.751 & 66 & 0.799 & 12 \\
before\_attention\_norm\_blocks\_9\_16\_token\_mask & 80 & 0.916 & 66 & 0.948 & 1 \\
before\_mlp & 46 & 0.751 & 46 & 0.679 & 16 \\
pre\_residual\_blocks\_17\_24 & 82 & 0.917 & 45 & 0.960 & 0 \\
before\_mlp\_residual & 42 & 0.817 & 42 & 0.884 & 1 \\
after\_embedding\_token\_mask & 33 & 0.666 & 33 & 0.672 & 10 \\
before\_attention\_norm\_blocks\_17\_24\_token\_mask & 80 & 0.961 & 23 & 0.976 & 0 \\
after\_mlp\_residual\_token\_mask & 22 & 0.831 & 22 & 0.758 & 5 \\
before\_attention\_channel\_mask & 22 & 0.735 & 22 & 0.780 & 4 \\
before\_mlp\_token\_mask & 22 & 0.711 & 22 & 0.707 & 7 \\
attention\_norm\_out\_gain\_scale & 11 & 0.796 & 11 & 0.901 & 0 \\
\bottomrule
\end{tabular}
\end{adjustbox}
\end{table}

% generated by scripts/paper/tab_swin.py from results/swin_*/psbd_metrics.json (90 implanted cells), checkpoints/swin_*/args.json at 2026-09-23T16:19:12.040053+00:00, commit afe5fd94d50b77f30f654fe075b768040987adef-dirty. Do not edit by hand.
\begin{table}[htbp]
\centering
\small
\caption{Swin-S paired placement gains at the headline quantile over the models carrying both placements, 5000-resample bootstrap intervals.}
\label{tab:swin-gains}
\begin{adjustbox}{max width=\linewidth}
\begin{tabular}{llrrl}
\toprule
comparison & rule & n & mean delta AUROC & 95\% CI \\
\midrule
token mask, attention input minus dropout, after both residual adds & matched & 80 & +0.165 & [+0.116, +0.217] \\
token mask, attention input minus dropout, after both residual adds & adaptive & 80 & +0.115 & [+0.070, +0.162] \\
dropout, attention input minus dropout, after both residual adds & matched & 80 & +0.026 & [$-$0.022, +0.076] \\
dropout, attention input minus dropout, after both residual adds & adaptive & 80 & +0.020 & [$-$0.037, +0.079] \\
token mask, attention input minus dropout, before both residual adds & matched & 80 & +0.148 & [+0.107, +0.191] \\
token mask, attention input minus dropout, before both residual adds & adaptive & 80 & +0.124 & [+0.077, +0.174] \\
token mask, attention input minus token mask, attention output before the add & matched & 80 & +0.157 & [+0.111, +0.204] \\
token mask, attention input minus token mask, attention output before the add & adaptive & 80 & +0.136 & [+0.086, +0.187] \\
token mask, attention input minus dropout, attention input & matched & 80 & +0.138 & [+0.097, +0.184] \\
token mask, attention input minus dropout, attention input & adaptive & 80 & +0.095 & [+0.057, +0.137] \\
\bottomrule
\end{tabular}
\end{adjustbox}
\end{table}

